\documentclass[aps,prd,reprint,preprintnumbers,showpacs,floatfix,nofootinbib,superscript address,longbibliography]{revtex4-1}
\pdfoutput=1

\usepackage[utf8]{inputenc}
\usepackage{scalerel,stackengine}%	omit indices
\usepackage{titlesec}%	PRL paragraph	
\usepackage{parskip}
\usepackage{amssymb}
\usepackage{stix}
\usepackage{hhline}
\usepackage{amsmath}
\usepackage{mathtools}
\usepackage[dvipsnames]{xcolor}
\usepackage{xspace}
\usepackage{multirow,tabularx}
\usepackage{siunitx}
\usepackage{multirow}
\usepackage{graphicx}
\usepackage{xstring}
\usepackage{etoolbox}
\usepackage{notoccite}
\usepackage{natbib}
\usepackage{mathrsfs}
\usepackage{lineno}
\usepackage{tensor}
\usepackage{accents}
\usepackage{tikz}

\allowdisplaybreaks

\parskip 1mm
\parindent 2mm

%%=============%%
%%  SZ macros  %%
%%=============%%

\newcommand*{\ie}{i.e., }
\newcommand*{\eg}{e.g., }
\newcommand*{\fig}{Fig.\@\xspace}
\newcommand*{\eq}{eq.\@\xspace}
\newcommand*{\eqs}{eqs.\@\xspace}
\newcommand*{\wrt}{w.r.t.\@\xspace}
\newcommand*{\dof}{d.o.f.\@\xspace}

%	Comments by SZ
%\newcommand{\sz}[1]{{\color[RGB]{0,34,170}{#1}}}
%\newcommand{\wb}[1]{{\color[RGB]{0,100,154}{#1}}}
%\newcommand{\szC}[1]{{\color[RGB]{0,34,170}{[\textbf{SZ: #1]}}}}

%%=============%%
%%  WB macros  %%
%%=============%%

%	PRL paragraph
\newrobustcmd{\pea}[1]{%
	\emph{#1}\textbf{\ \ \ ---}
}
\titleformat{\paragraph}[runin]{\normalfont\normalsize\bfseries}{\emph\theparagraph}{1em}{\pea}

%	The kinetic coefficient, with a nice, conservative symbol
\newrobustcmd{\alp}[1]{%
       {\alpha}	
}

%	The mass `precursor', also with a nice obvious symbol now
\newrobustcmd{\bet}[1]{%
  \tensor[^{(#1)}]{\mu}{}
}

%	ommission of indices
\stackMath 
\newcommand\OmitIndices[1]{% 
\savestack{\tmpbox}{\stretchto{%   
\scaleto{%     
\scalerel*[\widthof{\ensuremath{#1}}]{\kern-.6pt\curlywedge\kern-.6pt}%     
{\rule[-\textheight/2]{1ex}{\textheight}}%WIDTH-LIMITED BIG WEDGE   
}{\textheight}%  
}{0.5ex}}% 
\stackon[1pt]{#1}{\tmpbox}% 
}

%	macros for non-Riemannian geometry
\newrobustcmd{\LPV}[1]{%
	%{L_{\text{PV}}}%
	\IfEqCase{#1}{%
	{}{L_{\text{PV}}}%
	}%
	[{L_{\text{PV}}\left(#1\right)}]%
}


\newrobustcmd{\F}[2][placeholder]{%
	\IfEqCase{#1}{%
	{placeholder}{\tensor{T}{#2}}%
	{1}{\tensor[^{(1)}]{F}{#2}}%
	{2}{\tensor[^{(2)}]{F}{#2}}%
	{3}{\tensor[^{(3)}]{F}{#2}}%
	}%
	[\packageError{cosmicclass}{Symbol #1 is not an irreducible part!}{}]%
}


%	macros for non-Riemannian geometry

\newrobustcmd{\GenericVector}[1]{%
	\smash{{#1}_{\tensor[^{{(2)}}]{\hspace{-1pt}\lambda}{}}^{J^P}}%
}
\newrobustcmd{\GenericTensor}[1]{%
	\smash{{#1}_{\tensor[^{{(1)}}]{\hspace{-1pt}\lambda}{}}^{J^P}}%
}

\newrobustcmd{\qe}[1]{%
	\tensor{q}{_e}%
}
\newrobustcmd{\dqe}[1]{%
	\tensor{\dot{q}}{_e}%
}
\newrobustcmd{\qw}[1]{%
	\tensor{q}{_\omega}%
}
\newrobustcmd{\dqw}[1]{%
	\tensor{\dot{q}}{_\omega}%
}




\newrobustcmd{\chie}[1]{%
	\tensor{\chi}{_e}%
}
\newrobustcmd{\chiw}[1]{%
	\tensor{\chi}{_\omega}%
}
\newrobustcmd{\phil}[1]{%
	\tensor{\phi}{_\lambda}%
}
\newrobustcmd{\phie}[1]{%
	\tensor{\phi}{_e}%
}
\newrobustcmd{\phiw}[1]{%
	\tensor{\phi}{_\omega}%
}
\newrobustcmd{\pil}[1]{%
	\tensor{p}{_\lambda}%
}
\newrobustcmd{\pie}[1]{%
	\tensor{p}{_e}%
}
\newrobustcmd{\piw}[1]{%
	\tensor{p}{_\omega}%
}



\newrobustcmd{\g}[1]{%
	\tensor{g}{#1}%
}
\newrobustcmd{\rcCon}[1]{%
	\tensor{\Gamma}{#1}%
}
\newrobustcmd{\rCon}[1]{%
	\tensor*{\mathring{\Gamma}}{#1}%
}
\newrobustcmd{\B}[1]{%
	\tensor{B}{#1}%
}
\newrobustcmd{\PD}[1]{%
	\tensor{\partial}{#1}%
}
\newrobustcmd{\rD}[1]{%
	\tensor{\mathring{\nabla}}{#1}%
}
\newrobustcmd{\rcD}[1]{%
	\tensor{\nabla}{#1}%
}
\newrobustcmd{\rR}[1]{%
	\tensor{\mathring{R}}{#1}%
}
\newrobustcmd{\rcR}[1]{%
	\tensor{R}{#1}%
}

%I had to change this from "newrobustcmd" to "renewrobustcmd" to avoid an error
\renewrobustcmd{\F}[2][placeholder]{%
	\IfEqCase{#1}{%
	{placeholder}{\tensor{T}{#2}}%
	{1}{\tensor[^{(1)}]{F}{#2}}%
	{2}{\tensor[^{(2)}]{F}{#2}}%
	{3}{\tensor[^{(3)}]{F}{#2}}%
	}%
	[\packageError{cosmicclass}{Symbol #1 is not an irreducible part!}{}]%
}

\newrobustcmd{\Si}[2][placeholder]{%
	\IfEqCase{#1}{%
	{placeholder}{\tensor{T}{#2}}%
	{1}{\tensor[^{(1)}]{S}{#2}}%
	{2}{\tensor[^{(2)}]{S}{#2}}%
	{3}{\tensor[^{(3)}]{S}{#2}}%
	}%
	[\packageError{cosmicclass}{Symbol #1 is not an irreducible part!}{}]%
}

\newrobustcmd{\T}[2][placeholder]{%
	\IfEqCase{#1}{%
	{placeholder}{\tensor{T}{#2}}%
	{1}{\tensor[^{(1)}]{T}{#2}}%
	{2}{\tensor[^{(2)}]{T}{#2}}%
	{3}{\tensor[^{(3)}]{T}{#2}}%
	}%
	[\packageError{cosmicclass}{Symbol #1 is not an irreducible part!}{}]%
}

\newrobustcmd{\mass}[2][placeholder]{%
	\IfEqCase{#1}{%
	{placeholder}{\tensor{m}{#2}}%
	{1}{\tensor[^{(1)}]{m}{#2}}%
	{2}{\tensor[^{(2)}]{m}{#2}}%
	{3}{\tensor[^{(3)}]{m}{#2}}%
	}%
	[\packageError{cosmicclass}{Symbol #1 is not an irreducible part!}{}]%
}

\newrobustcmd{\TLambda}[2][placeholder]{%
	\IfEqCase{#1}{%
	{placeholder}{\tensor{\lambda}{#2}}%
	{1}{\tensor[^{(1)}]{\lambda}{#2}}%
	{2}{\tensor[^{(2)}]{\lambda}{#2}}%
	{3}{\tensor[^{(3)}]{\lambda}{#2}}%
	{Q}{\tensor[^{(Q)}]{\lambda}{#2}}%
	}%
	[\packageError{cosmicclass}{Symbol #1 is not an irreducible part!}{}]%
}

\newrobustcmd{\RLambda}[2][placeholder]{%
	\IfEqCase{#1}{%
	{placeholder}{\tensor{\lambda}{#2}}%
	{1}{\tensor[^1]{\lambda}{#2}}%
	{2}{\tensor[^2]{\lambda}{#2}}%
	{3}{\tensor[^3]{\lambda}{#2}}%
	{R}{\tensor[^{(R)}]{\lambda}{#2}}%
	}%
	[\packageError{cosmicclass}{Symbol #1 is not an irreducible part!}{}]%
}

\newrobustcmd{\QLambda}[2][placeholder]{%
	\IfEqCase{#1}{%
		{placeholder}{\tensor{\hat{\lambda}}{#2}}%
	{1}{\tensor[^1]{\hat{\lambda}}{#2}}%
	{2}{\tensor[^2]{\hat{\lambda}}{#2}}%
	{3}{\tensor[^3]{\hat{\lambda}}{#2}}%
	}%
	[\packageError{cosmicclass}{Symbol #1 is not an irreducible part!}{}]%
}





\iffalse

\newrobustcmd{\PiP}[2][placeholder]{%
\IfEqCase{#2}{%
{B0p}{\hat{\pi}}%
{B1p}{\tensor{\overset{\wedge}{\hat{\pi}}}{_{\indiq[#1]{2}}}}%
{B1m}{\tensor{\hat{\pi}}{_{\perp\indiq[#1]{1}}}}%
{B2p}{\tensor{\overset{\sim}{\hat{\pi}}}{_{\indiq[#1]{2}}}}%
{A0p}{\tensor{\hat{\pi}}{_\perp}}%
{A0m}{\tensor[^{\text{P}}]{\hat{\pi}}{}}%
{A1p}{\tensor{\overset{\wedge}{\hat{\pi}}}{_{\perp\indiq[#1]{2}}}}%
{A1m}{\tensor{\overset{\rightharpoonup}{\hat{\pi}}}{_{\indiq[#1]{1}}}}%
{A2p}{\tensor{\overset{\sim}{\hat{\pi}}}{_{\perp\indiq[#1]{2}}}}%
{A2m}{\tensor[^{\text{T}}]{\hat{\pi}}{_{\indiq[#1]{3}}}}%
}[\packageError{cosmicclass}{Unidentified Critical Case: #1}{}]%
}

\fi


\newrobustcmd{\Mgra}[1]{%
  {\tensor{M}{_{\text{#1}}}}%
}
\newrobustcmd{\Mpro}[1]{%
  {\tensor{\mathproper{M}}{_{\text{#1}}}}%
}
\newrobustcmd{\Malt}[1]{%
  {\tensor{\mathscr{M}}{_{\text{#1}}}}%
}
\newrobustcmd{\Mkom}[1]{%
  {\tensor{\mathfrak{M}}{_{\text{#1}}}}%
}

\newrobustcmd{\Mtotal}{%
  {\tensor{M}{_{\text{T}}}}%
}

\newrobustcmd{\Qtotal}{%
  {\tensor{Q}{_{\text{T}}}}%
}

\newrobustcmd{\Qtotalcal}{%
  {\tensor{\mathcal{  Q}}{_{\text{T}}}}%
}

\newrobustcmd{\action}[1]{%
  {\tensor{S}{_{\text{#1}}}}%
}

\newrobustcmd{\lagrangian}[1]{%
  {\tensor{L}{_{\text{#1}}}}%
}

\newrobustcmd{\lagrangianprop}[1]{%
  {\tensor{\mathproper{L}}{_{\text{#1}}}}%
}

%\newcommand{\ifeq}{\mathrel{\smash[t]{\stackrel{{\text{(e.g.)}}}{=}}}} %analogue equality
%\newcommand{\aneq}{\mathrel{\smash[t]{\stackrel{{\text{an}}}{=}}}} %analogue equality

\newrobustcmd{\epl}{%
  {\tensor{\mathsf{e}}{_+}}%
}

\newrobustcmd{\epe}{%
  {\tensor{\mathsf{e}}{_\perp}}%
}

%\newrobustcmd{\mul}[1]{%
%  {\text{mul}(#1)}%
%}

%cell color
\newrobustcmd{\qz}{%
  {\text{\color{orange}\cmark}}%
}
\newrobustcmd{\jz}{%
  {\text{\color{red}\xmark}}%
}

\newrobustcmd{\projmatrix}[2][placeholder]{%
  {\tensor*{M}{_{#1}^{#2}}}
}
\newrobustcmd{\projorthhum}[2][placeholder]{%
  {\tensor[^#2]{\smash{\check{\mathcal{  P}}}}{#1}}
}
\newrobustcmd{\projorthhumu}[2][placeholder]{%
  {\tensor[^#2]{\smash{\check{\mathcal{  P}}}}{#1}}
}
\newrobustcmd{\projorth}[2][placeholder]{%
  {\tensor[^#2]{\smash{\hat{\mathcal{  P}}}}{#1}}
}
\newrobustcmd{\projlore}[2][placeholder]{%
  {\tensor[^#2]{\hat{\mathcal{  P}}}{#1}}
}


\newrobustcmd{\gensec}[3][placeholder]{%
  {\tensor*[^#1]{\chi}{^{#2}_{\acu{#3}}}}
}

%general linear group in four real dimensions
\newrobustcmd{\glfourr}{%
  {\mathrm{GL}(4,\mathbb{R})}%
}

%special linear group in two complex dimensions
\newrobustcmd{\sltwoc}{%
  {\mathrm{SL}(2,\mathbb{C})}%
}

%poincare group
\newrobustcmd{\poincare}{%
  {\mathbb{R}^{1,3}\rtimes\mathrm{SO}^+(1,3)}%
}

%poincare group formal
\newrobustcmd{\poincaref}{%
  {\mathrm{P}(1,3)}%
}

%weyl group
\newrobustcmd{\weyl}{%
  {\mathrm{W}(1,3)}%
}

%weyl group
\newrobustcmd{\conformal}{%
  {\mathrm{C}(1,3)}%
}

%diffeomorphism group
\newrobustcmd{\diffeomorphism}{%
  {\mathbb{R}^{1,3}}%
}

%lorentz group
\newrobustcmd{\soonethree}{%
  {\mathrm{SO}^+(1,3)}%
}

%special orthogonal group
\newrobustcmd{\othree}{%
  {\mathrm{SO}(3)}%
}

%special orthogonal group
\newrobustcmd{\sothree}{%
  {\mathrm{SO}(3)}%
}

%special orthogonal group
\newrobustcmd{\sotwo}{%
  {\mathrm{SO}(2)}%
}

%QCD group
\newrobustcmd{\suthreec}{%
  {\mathrm{SU}(3)_{\text{c}}}%
}

%weak isospin group
\newrobustcmd{\sutwol}{%
  {\mathrm{SU}(2)_{\text{L}}}%
}

%weak hypercharge group
\newrobustcmd{\uoney}{%
  {\mathrm{U}(1)_{\text{Y}}}%
}

%weak hypercharge group
\newrobustcmd{\uone}{%
  {\mathrm{U}(1)}%
}

%electromagnetic group
\newrobustcmd{\uoneem}{%
  {\mathrm{U}(1)_{\text{em}}}%
}

\newrobustcmd{\sutwo}{%
  {\mathrm{SU}(2)}%
}

\newrobustcmd{\eplus}{%
  {\tensor{\mathsf{e}}{_{+}}}%
}

\newrobustcmd{\esf}[1]{%
  {\tensor{\mathsf{e}}{_{#1}}}
}%
\newrobustcmd{\esfu}[1]{%
  {\tensor{\mathsf{e}}{^{#1}}}
}%
\newrobustcmd{\gam}[1]{%
  {\tensor{\gamma}{_{#1}}}
}%
\newrobustcmd{\gamu}[1]{%
  {\tensor{\gamma}{^{#1}}}
}%

\newrobustcmd{\planck}{%
  {m_{\text{p}}}%
}

\newrobustcmd{\Planck}{%
	{M_{\text{Pl}}}%
}

\newrobustcmd{\caligR}{%
  {\mathcal{R}}%
}
\newrobustcmd{\caligT}{%
  {\mathcal{T}}%
}

\newrobustcmd{\pgt}{%
  PGT\textsuperscript{q,+}\ %
}


\newrobustcmd{\unl}[1]{%
  {\mathfrak{#1}}%
}
\newrobustcmd{\ovl}[1]{%
\overline{#1}%
}
\newrobustcmd{\acu}[1]{%
\acute{#1}%
}


\newrobustcmd{\indiq}[2][placeholder]{%
\IfEqCase{#1}{%
  {placeholder}{%
    \IfEqCase{#2}{%
      {1}{\ovl{k}}%
      {2}{\ovl{kl}}%
      {3}{\ovl{klm}}%
    }%
  }%
}[#1]%
}%


\newrobustcmd{\indaq}[2][placeholder]{%
\IfEqCase{#1}{%
  {placeholder}{%
    \IfEqCase{#2}{%
      {1}{\overline{k}}%
      {2}{\overline{kl}}%
      {3}{\overline{klm}}%
    }%
  }%
}[#1]%
}%

\newrobustcmd{\indeq}[2][placeholder]{%
\IfEqCase{#1}{%
  {placeholder}{%
    \IfEqCase{#2}{%
      {1}{k}%
      {2}{kl}%
      {3}{klm}%
    }%
  }%
}[#1]%
}%

\newrobustcmd{\indoq}[2][placeholder]{%
\IfEqCase{#1}{%
  {placeholder}{%
    \IfEqCase{#2}{%
      {1}{\alpha}%
      {2}{\alpha\beta}%
      {3}{\alpha\beta\gamma}%
    }%
  }%
}[#1]%
}%

\newrobustcmd{\fcphi}[1]{%
  \tensor[^{\text{FC}}]{\phi}{_{#1}}%
}

\newrobustcmd{\scphi}[1]{%
  \tensor[^{\text{SC}}]{\phi}{_{#1}}%
}

\newrobustcmd{\fcmul}[1]{%
  \tensor[^{\text{FC}}]{\upsilon}{_{#1}}%
}

\newrobustcmd{\arb}{%
  {\tensor{f}{_{\text{lin}}}}%
}

\newrobustcmd{\scmul}[1]{%
  \tensor[^{\text{SC}}]{\upsilon}{_{#1}}%
}

\newrobustcmd{\foli}[1]{%
\tensor{n}{_{#1}}%
}

\newrobustcmd{\foliu}[1]{%
\tensor{n}{^{#1}}%
}

\newrobustcmd{\covderl}[1]{%
\tensor{\mathcal{D}}{^{\flat}_{\indiq[#1]{1}}}%
}

\newrobustcmd{\covder}[1]{%
\tensor{\mathcal{D}}{_{\indiq[#1]{1}}}%
}
\newrobustcmd{\coder}[1]{%
\tensor{D}{_{\indiq[#1]{1}}}%
}

\newrobustcmd{\deltal}[2]{%
  \tensor*{\delta}{_{\phantom{\flat}}^{\flat}_{#1}^{#2}}%
}


\newrobustcmd{\deltaud}[2]{%
  \tensor*{\delta}{^{#1}_{#2}}%
}
\newrobustcmd{\etau}[1]{%
\tensor{\eta}{^{\indiq[#1]{2}}}%
}
\newrobustcmd{\etaul}[1]{%
\tensor{\eta}{^{\flat}^{\indiq[#1]{2}}}%
}
\newrobustcmd{\etad}[1]{%
\tensor{\eta}{_{\indiq[#1]{2}}}%
}
\newrobustcmd{\etadl}[1]{%
\tensor{\eta}{^{\flat}_{\indiq[#1]{2}}}%
}

\newrobustcmd{\epsul}[1]{%
\tensor{\epsilon}{^{\flat}^{\indiq[#1]{3}}^{\perp}}
}
\newrobustcmd{\epsdl}[1]{%
\tensor{\epsilon}{^{\flat}_{\indiq[#1]{3}}_{\perp}}
}
\newrobustcmd{\epsd}[1]{%
\tensor{\epsilon}{_{\indiq[#1]{3}}_{\perp}}
}
\newrobustcmd{\epsu}[1]{%
\tensor{\epsilon}{^{\indiq[#1]{3}}^{\perp}}
}

\newrobustcmd{\hfl}[2]{%
  \tensor{h}{^{\flat}_{#1}^{#2}}
}

\newrobustcmd{\cgalp}{\tensor{\alpha}{_{\text{CG}}}}

\newrobustcmd{\cbet}[1]{%
  \tensor{\bar{\beta}}{_{#1}}
}
\newrobustcmd{\calp}[1]{%
  \tensor{\bar{\alpha}}{_{#1}}
}


\newrobustcmd{\alpg}[1]{%
  \tensor{\check{\alpha}}{_{#1}}
}
\newrobustcmd{\betg}[1]{%
  \tensor{\check{\beta}}{_{#1}}
}
\newrobustcmd{\calpg}[1]{%
  \tensor{\acu{\alpha}}{_{#1}}
}
\newrobustcmd{\cbetg}[1]{%
  \tensor{\acu{\beta}}{_{#1}}
}

\newrobustcmd{\hub}{%
  {\underline{\mathsf{h}}}
}
\newrobustcmd{\hubm}{%
  {\underline{\mathsf{h}}^{-1}}
}
\newrobustcmd{\hob}{%
  {\bar{\mathsf{h}}}
}
\newrobustcmd{\hobm}{%
  {\bar{\mathsf{h}}^{-1}}
}
\newrobustcmd{\hdet}{%
  {\det \mathsf{h}}
}
\newrobustcmd{\hmdet}{%
  {\det \mathsf{h}^{-1}}
}
\newrobustcmd{\Rsf}{%
  {\mathsf{R}}
}


\newrobustcmd{\alpm}[2][placeholder]{%
  {\tensor*{\hat{\alpha}}{_{#1}^{#2}}}
}
\newrobustcmd{\calpm}[2][placeholder]{%
  \tensor*{\bar{\alpha}}{_{#1}^{#2}}
}

\newrobustcmd{\betm}[2][placeholder]{%
  {\tensor*{\hat{\beta}}{_{#1}^{#2}}}
}
\newrobustcmd{\cbetm}[2][placeholder]{%
  \tensor*{\bar{\beta}}{_{#1}^{#2}}
}

\newrobustcmd{\lamr}{%
  {\zeta_{\mathcal{  R}} }
}
\newrobustcmd{\barlamr}{%
  {\bar{\zeta}_{\mathcal{  R}} }
}
\newrobustcmd{\lamt}{%
  {\zeta_{\mathcal{  T}} }
}
\newrobustcmd{\barlamt}{%
  {\bar{\zeta}_{\mathcal{  T}} }
}


%these are old, will need to delete
\newrobustcmd{\atmp}[1]{%
  \tensor{\hat{a}}{_{#1}}
}
\newrobustcmd{\btmp}[1]{%
  \tensor{b}{_{#1}}
}
%\newrobustcmd{\ctmp}[1]{%
%  \tensor{c}{_{#1}}
%}

\newrobustcmd{\ctmp}[2][placeholder]{%
  {\tensor*{c}{_{#1}^{#2}}}
}
\newrobustcmd{\dtmp}[2][placeholder]{%
  {\tensor*{d}{_{#1}^{#2}}}
}

\newrobustcmd{\etmp}[1]{%
  \tensor{e}{_{#1}}
}


\newrobustcmd{\batmp}[1]{%
  \tensor{\ovl{a}}{_{#1}}
}
\newrobustcmd{\bbtmp}[1]{%
  \tensor{\ovl{b}}{_{#1}}
}
\newrobustcmd{\bctmp}[1]{%
  \tensor{\ovl{c}}{_{#1}}
}
\newrobustcmd{\bdtmp}[1]{%
  \tensor{\ovl{d}}{_{#1}}
}
\newrobustcmd{\betmp}[1]{%
  \tensor{\ovl{e}}{_{#1}}
}

% SO(1,3) irreps of the fields

\newrobustcmd{\ptl}[1]{%
  \tensor{\partial}{#1}
}

\newrobustcmd{\etaf}[1]{%
  \tensor{\eta}{#1}
}

\newrobustcmd{\epsf}[1]{%
  \tensor{\epsilon}{#1}
}

\newrobustcmd{\RSO}[2][placeholder]{%
  {\tensor[^{#2}]{\mathcal{  R}}{#1}}
}
\newrobustcmd{\TSO}[2][placeholder]{%
  {\tensor[^{#2}]{\mathcal{  T}}{#1}}
}
\newrobustcmd{\FSO}[2][placeholder]{%
  {\tensor[^{#2}]{\mathcal{  F}}{#1}}
}
\newrobustcmd{\spinSO}[2][placeholder]{%
  {\tensor[^{#2}]{\sigma}{#1}}
}
\newrobustcmd{\RLambdaSO}[2][placeholder]{%
  {\tensor[^{#2}]{\zeta}{#1}}
}
\newrobustcmd{\TLambdaSO}[2][placeholder]{%
  {\tensor[^{#2}]{\zeta}{#1}}
}
\newrobustcmd{\KSO}[2][placeholder]{%
  {\tensor[^{#2}]{\mathcal{  K}}{#1}}
}

% linear gravity

\newrobustcmd{\bper}[2][placeholder]{%
\IfEqCase{#2}{%
  {s}{\tensor{\mathfrak{s}}{#1}}%
  {a}{\tensor{\mathfrak{a}}{#1}}%
  {sbar}{\tensor{\bar{\mathfrak{s}}}{#1}}%
}[\packageError{cosmicclass}{Unidentified Critical Case: #1}{}]%
}


% formal linear gravity (i.e. using the `flat' symbol)

\newrobustcmd{\Jl}{%
  {J^{\flat}}%
}%

\newrobustcmd{\Nl}{%
  {N^{\flat}}%
}%

\newrobustcmd{\haml}[2][placeholder]{%
\IfEqCase{#2}{%
{mom0p}{\tensor{\mathcal{H}}{^{\flat}_{\perp}}}%
{mom1m}{\tensor{\mathcal{H}}{^{\flat}_{\indoq[#1]{1}}}}%
{rot1p}{\tensor{\mathcal{H}}{^{\flat}_{\indaq[#1]{2}}}}%
{rot1m}{\tensor{\mathcal{H}}{^{\flat}_{\perp}_{\indaq[#1]{1}}}}%
}[\packageError{cosmicclass}{Unidentified Critical Case: #1}{}]%
}

\newrobustcmd{\arc}[2][placeholder]{%
\IfEqCase{#2}{%
{B1p}{\tensor{\vartheta}{_{\perp\indiq[#1]{2}}}}%
{B2m}{\tensor[^{\text{T}}]{\vartheta}{_{\indiq[#1]{3}}}}%
{A0m}{\tensor[^{\text{P}}]{\vartheta}{}}%
{A1p}{\tensor{\overset{\wedge}{\vartheta}}{_{\perp\indiq[#1]{2}}}}%
{A1m}{\tensor{\overset{\rightharpoonup}{\vartheta}}{_{\indiq[#1]{1}}}}%
{A2p}{\tensor{\overset{\sim}{\vartheta}}{_{\perp\indiq[#1]{2}}}}%
{A2m}{\tensor[^{\text{T}}]{\vartheta}{_{\perp\indiq[#1]{3}}}}%
}[\packageError{cosmicclass}{Unidentified Critical Case: #1}{}]%
}

\newrobustcmd{\pic}[2][placeholder]{%
\IfEqCase{#2}{%
{B0p}{\varphi}%
{B1p}{\tensor{\overset{\wedge}{\varphi}}{_{\indiq[#1]{2}}}}%
{B1m}{\tensor{\varphi}{_{\perp\indiq[#1]{1}}}}%
{B2p}{\tensor{\overset{\sim}{\varphi}}{_{\indiq[#1]{2}}}}%
{A0p}{\tensor{\varphi}{_\perp}}%
{A0m}{\tensor[^{\text{P}}]{\varphi}{}}%
{A1p}{\tensor{\overset{\wedge}{\varphi}}{_{\perp\indiq[#1]{2}}}}%
{A1m}{\tensor{\overset{\rightharpoonup}{\varphi}}{_{\indiq[#1]{1}}}}%
{A2p}{\tensor{\overset{\sim}{\varphi}}{_{\perp\indiq[#1]{2}}}}%
{A2m}{\tensor[^{\text{T}}]{\varphi}{_{\indiq[#1]{3}}}}%
}[\packageError{cosmicclass}{Unidentified Critical Case: #1}{}]%
}

\newrobustcmd{\picu}[2][placeholder]{%
\IfEqCase{#2}{%
{B0p}{\varphi}%
{B1p}{\tensor{\smash{\overset{\wedge}{\varphi}}}{^{\indiq[#1]{2}}}}%
{B1m}{\tensor{\varphi}{^{\perp\indiq[#1]{1}}}}%
{B2p}{\tensor{\smash{\overset{\sim}{\varphi}}}{^{\indiq[#1]{2}}}}%
{A0p}{\tensor{\varphi}{_\perp}}%
{A0m}{\tensor[^{\text{P}}]{\varphi}{}}%
{A1p}{\tensor{\smash{\overset{\wedge}{\varphi}}}{^{\perp\indiq[#1]{2}}}}%
{A1m}{\tensor{\smash{\overset{\rightharpoonup}{\varphi}}}{^{\indiq[#1]{1}}}}%
{A2p}{\tensor{\smash{\overset{\sim}{\varphi}}}{^{\perp\indiq[#1]{2}}}}%
{A2m}{\tensor[^{\text{T}}]{\varphi}{^{\indiq[#1]{3}}}}%
}[\packageError{cosmicclass}{Unidentified Critical Case: #1}{}]%
}

\newrobustcmd{\picl}[2][placeholder]{%
\IfEqCase{#2}{%
{B0p}{\tensor{\varphi}{^{\flat}}}%
{B1p}{\tensor{\smash{\overset{\wedge}{\varphi}}}{^{\flat}_{\indiq[#1]{2}}}}%
{B1m}{\tensor{\varphi}{^{\flat}_{\perp}_{\indiq[#1]{1}}}}%
{B2p}{\tensor{\smash{\overset{\sim}{\varphi}}}{^{\flat}_{\indiq[#1]{2}}}}%
{A0p}{\tensor{\varphi}{_\perp}^{\flat}}%
{A0m}{\tensor[^{\text{P}}]{\varphi}{^{\flat}}}%
{A1p}{\tensor{\smash{\overset{\wedge}{\varphi}}}{^{\flat}_{\perp\indiq[#1]{2}}}}%
{A1m}{\tensor{\smash{\overset{\rightharpoonup}{\varphi}}}{^{\flat}_{\indiq[#1]{1}}}}%
{A2p}{\tensor{\smash{\overset{\sim}{\varphi}}}{^{\flat}_{\perp\indiq[#1]{2}}}}%
{A2m}{\tensor[^{\text{T}}]{\varphi}{^{\flat}_{\indiq[#1]{3}}}}%
}[\packageError{cosmicclass}{Unidentified Critical Case: #1}{}]%
}

\newrobustcmd{\mull}[2][placeholder]{%
\IfEqCase{#2}{%
{B0p}{\tensor{u}{^{\flat}}}%
{B1p}{\tensor{\smash{\overset{\wedge}{u}}}{^{\flat}_{\indiq[#1]{2}}}}%
{B1m}{\tensor{u}{^{\flat}_{\perp}_{\indiq[#1]{1}}}}%
{B2p}{\tensor{\smash{\overset{\sim}{u}}}{^{\flat}_{\indiq[#1]{2}}}}%
{A0p}{\tensor{u}{_\perp}^{\flat}}%
{A0m}{\tensor[^{\text{P}}]{u}{^{\flat}}}%
{A1p}{\tensor{\smash{\overset{\wedge}{u}}}{^{\flat}_{\perp\indiq[#1]{2}}}}%
{A1m}{\tensor{\smash{\overset{\rightharpoonup}{u}}}{^{\flat}_{\indiq[#1]{1}}}}%
{A2p}{\tensor{\smash{\overset{\sim}{u}}}{^{\flat}_{\perp\indiq[#1]{2}}}}%
{A2m}{\tensor[^{\text{T}}]{u}{^{\flat}_{\indiq[#1]{3}}}}%
}[\packageError{cosmicclass}{Unidentified Critical Case: #1}{}]%
}

\newrobustcmd{\mul}[2][placeholder]{%
\IfEqCase{#2}{%
{B0p}{\tensor{u}{}}%
{B1p}{\tensor{\smash{\overset{\wedge}{u}}}{_{\indiq[#1]{2}}}}%
{B1m}{\tensor{u}{_{\perp}_{\indiq[#1]{1}}}}%
{B2p}{\tensor{\smash{\overset{\sim}{u}}}{_{\indiq[#1]{2}}}}%
{A0p}{\tensor{u}{_\perp}}%
{A0m}{\tensor[^{\text{P}}]{u}{}}%
{A1p}{\tensor{\smash{\overset{\wedge}{u}}}{_{\perp\indiq[#1]{2}}}}%
{A1m}{\tensor{\smash{\overset{\rightharpoonup}{u}}}{_{\indiq[#1]{1}}}}%
{A2p}{\tensor{\smash{\overset{\sim}{u}}}{_{\perp\indiq[#1]{2}}}}%
{A2m}{\tensor[^{\text{T}}]{u}{_{\indiq[#1]{3}}}}%
}[\packageError{cosmicclass}{Unidentified Critical Case: #1}{}]%
}

\newrobustcmd{\PiP}[2][placeholder]{%
\IfEqCase{#2}{%
{B0p}{\hat{\pi}}%
{B1p}{\tensor{\overset{\wedge}{\hat{\pi}}}{_{\indiq[#1]{2}}}}%
{B1m}{\tensor{\hat{\pi}}{_{\perp\indiq[#1]{1}}}}%
{B2p}{\tensor{\overset{\sim}{\hat{\pi}}}{_{\indiq[#1]{2}}}}%
{A0p}{\tensor{\hat{\pi}}{_\perp}}%
{A0m}{\tensor[^{\text{P}}]{\hat{\pi}}{}}%
{A1p}{\tensor{\overset{\wedge}{\hat{\pi}}}{_{\perp\indiq[#1]{2}}}}%
{A1m}{\tensor{\overset{\rightharpoonup}{\hat{\pi}}}{_{\indiq[#1]{1}}}}%
{A2p}{\tensor{\overset{\sim}{\hat{\pi}}}{_{\perp\indiq[#1]{2}}}}%
{A2m}{\tensor[^{\text{T}}]{\hat{\pi}}{_{\indiq[#1]{3}}}}%
}[\packageError{cosmicclass}{Unidentified Critical Case: #1}{}]%
}

\newrobustcmd{\PiPu}[2][placeholder]{%
\IfEqCase{#2}{%
{B0p}{\hat{\pi}}%
{B1p}{\tensor{\smash{\overset{\wedge}{\hat{\pi}}}}{^{\indiq[#1]{2}}}}%
{B1m}{\tensor{\smash{\hat{\pi}}}{^{\perp\indiq[#1]{1}}}}%
{B2p}{\tensor{\smash{\overset{\sim}{\hat{\pi}}}}{^{\indiq[#1]{2}}}}%
{A0p}{\tensor{\smash{\hat{\pi}}}{^\perp}}%
{A0m}{\tensor[^{\text{P}}]{\smash{\hat{\pi}}}{}}%
{A1p}{\tensor{\smash{\overset{\wedge}{\hat{\pi}}}}{^{\perp\indiq[#1]{2}}}}%
{A1m}{\tensor{\smash{\overset{\rightharpoonup}{\hat{\pi}}}}{^{\indiq[#1]{1}}}}%
{A2p}{\tensor{\smash{\overset{\sim}{\hat{\pi}}}}{^{\perp\indiq[#1]{2}}}}%
{A2m}{\tensor[^{\text{T}}]{\smash{\hat{\pi}}}{^{\indiq[#1]{3}}}}%
}[\packageError{cosmicclass}{Unidentified Critical Case: #1}{}]%
}

\newrobustcmd{\sicl}[2][placeholder]{%
\IfEqCase{#2}{%
{B0p}{\tensor{\chi}{^{\flat}}}%
{B1p}{\tensor{\smash{\overset{\wedge}{\chi}}}{^{\flat}_{\indiq[#1]{2}}}}%
{B1m}{\tensor{\chi}{^{\flat}_{\perp}_{\indiq[#1]{1}}}}%
{B2p}{\tensor{\smash{\overset{\sim}{\chi}}}{^{\flat}_{\indiq[#1]{2}}}}%
{A0p}{\tensor{\chi}{^{\flat}_\perp}}%
{A0m}{\tensor[^{\text{P}}]{\chi}{^{\flat}}}%
{A1p}{\tensor{\smash{\overset{\wedge}{\chi}}}{^{\flat}_{\perp\indiq[#1]{2}}}}%
{A1m}{\tensor{\smash{\overset{\rightharpoonup}{\chi}}}{^{\flat}_{\indiq[#1]{1}}}}%
{A2p}{\tensor{\smash{\overset{\sim}{\chi}}}{^{\flat}_{\perp\indiq[#1]{2}}}}%
{A2m}{\tensor[^{\text{T}}]{\chi}{^{\flat}_{\indiq[#1]{3}}}}%
}[\packageError{cosmicclass}{Unidentified Critical Case: #1}{}]%
}

\newrobustcmd{\ticl}[2][placeholder]{%
\IfEqCase{#2}{%
{B0p}{\tensor{\zeta}{^{\flat}}}%
{B1p}{\tensor{\smash{\overset{\wedge}{\zeta}}}{^{\flat}_{\indiq[#1]{2}}}}%
{B1m}{\tensor{\zeta}{^{\flat}_{\perp}_{\indiq[#1]{1}}}}%
{B2p}{\tensor{\smash{\overset{\sim}{\zeta}}}{^{\flat}_{\indiq[#1]{2}}}}%
{A0p}{\tensor{\zeta}{^{\flat}_\perp}}%
{A0m}{\tensor[^{\text{P}}]{\zeta}{^{\flat}}}%
{A1p}{\tensor{\smash{\overset{\wedge}{\zeta}}}{^{\flat}_{\perp\indiq[#1]{2}}}}%
{A1m}{\tensor{\smash{\overset{\rightharpoonup}{\zeta}}}{^{\flat}_{\indiq[#1]{1}}}}%
{A2p}{\tensor{\smash{\overset{\sim}{\zeta}}}{^{\flat}_{\perp\indiq[#1]{2}}}}%
{A2m}{\tensor[^{\text{T}}]{\zeta}{^{\flat}_{\indiq[#1]{3}}}}%
}[\packageError{cosmicclass}{Unidentified Critical Case: #1}{}]%
}

\newrobustcmd{\PiPl}[2][placeholder]{%
\IfEqCase{#2}{%
{B0p}{\tensor{\hat{\pi}}{^{\flat}}}%
{B1p}{\tensor{\smash{\overset{\wedge}{\hat{\pi}}}}{^{\flat}_{\indiq[#1]{2}}}}%
{B1m}{\tensor{\hat{\pi}}{^{\flat}_{\perp}_{\indiq[#1]{1}}}}%
{B2p}{\tensor{\smash{\overset{\sim}{\hat{\pi}}}}{^{\flat}_{\indiq[#1]{2}}}}%
{A0p}{\tensor{\hat{\pi}}{_\perp}^{\flat}}%
{A0m}{\tensor[^{\text{P}}]{\hat{\pi}}{^{\flat}}}%
{A1p}{\tensor{\smash{\overset{\wedge}{\hat{\pi}}}}{^{\flat}_{\perp\indiq[#1]{2}}}}%
{A1m}{\tensor{\smash{\overset{\rightharpoonup}{\hat{\pi}}}}{^{\flat}_{\indiq[#1]{1}}}}%
{A2p}{\tensor{\smash{\overset{\sim}{\hat{\pi}}}}{^{\flat}_{\perp\indiq[#1]{2}}}}%
{A2m}{\tensor[^{\text{T}}]{\hat{\pi}}{^{\flat}_{\indiq[#1]{3}}}}%
}[\packageError{cosmicclass}{Unidentified Critical Case: #1}{}]%
}

\newrobustcmd{\sic}[2][placeholder]{%
\IfEqCase{#2}{%
{B0p}{\chi}%
{B1p}{\tensor{\overset{\wedge}{\chi}}{_{\indiq[#1]{2}}}}%
{B1m}{\tensor{\chi}{_{\perp\indiq[#1]{1}}}}%
{B2p}{\tensor{\overset{\sim}{\chi}}{_{\indiq[#1]{2}}}}%
{A0p}{\tensor{\chi}{_\perp}}%
{A0m}{\tensor[^{\text{P}}]{\chi}{}}%
{A1p}{\tensor{\overset{\wedge}{\chi}}{_{\perp\indiq[#1]{2}}}}%
{A1m}{\tensor{\overset{\rightharpoonup}{\chi}}{_{\indiq[#1]{1}}}}%
{A2p}{\tensor{\overset{\sim}{\chi}}{_{\perp\indiq[#1]{2}}}}%
{A2m}{\tensor[^{\text{T}}]{\chi}{_{\indiq[#1]{3}}}}%
}[\packageError{cosmicclass}{Unidentified Critical Case: #1}{}]%
}

\newrobustcmd{\lorsicpar}[2][placeholder]{%
\IfEqCase{#2}{%
{B0m}{\tensor*[^{\text{P}}]{\smash{\underline{\chi}}}{^{\parallel}}}%
{B1p}{\tensor*{\smash{\overset{\wedge}{\chi}}}{^{\parallel}_{\indiq[#1]{2}}}}%
{B1m}{\tensor*{\chi}{^{\parallel}_{\perp\indiq[#1]{1}}}}%
{B2m}{\tensor*[^{\text{T}}]{\smash{\underline{\chi}}}{^{\parallel}_{\indiq[#1]{3}}}}%
{A0p}{\tensor*{\chi}{^{\parallel}_\perp}}%
{A0m}{\tensor*[^{\text{P}}]{\chi}{^{\parallel}}}%
{A1p}{\tensor*{\smash{\overset{\wedge}{\chi}}}{^{\parallel}_{\perp\indiq[#1]{2}}}}%
{A1m}{\tensor*{\smash{\overset{\rightharpoonup}{\chi}}}{^{\parallel}_{\indiq[#1]{1}}}}%
{A2p}{\tensor*{\smash{\overset{\sim}{\chi}}}{^{\parallel}_{\perp\indiq[#1]{2}}}}%
{A2m}{\tensor*[^{\text{T}}]{\chi}{^{\parallel}_{\indiq[#1]{3}}}}%
}[\packageError{cosmicclass}{Unidentified Critical Case: #1}{}]%
}

\newrobustcmd{\lorsicpir}[2][placeholder]{%
\IfEqCase{#2}{%
{B0p}{\tensor*{\chi}{^{\vDash}}}%
{B1p}{\tensor*{\smash{\overset{\wedge}{\chi}}}{^{\vDash}_{\indiq[#1]{2}}}}%
{B1m}{\tensor*{\chi}{^{\vDash}_{\perp\indiq[#1]{1}}}}%
{B2p}{\tensor*{\smash{\overset{\sim}{\chi}}}{^{\vDash}_{\indiq[#1]{2}}}}%
{A0p}{\tensor*{\chi}{^{\vDash}_\perp}}%
{A0m}{\tensor*[^{\text{P}}]{\chi}{^{\vDash}}}%
{A1p}{\tensor*{\smash{\overset{\wedge}{\chi}}}{^{\vDash}_{\perp\indiq[#1]{2}}}}%
{A1m}{\tensor*{\smash{\overset{\rightharpoonup}{\chi}}}{^{\vDash}_{\indiq[#1]{1}}}}%
{A2p}{\tensor*{\smash{\overset{\sim}{\chi}}}{^{\vDash}_{\perp\indiq[#1]{2}}}}%
{A2m}{\tensor*[^{\text{T}}]{\chi}{^{\vDash}_{\indiq[#1]{3}}}}%
}[\packageError{cosmicclass}{Unidentified Critical Case: #1}{}]%
}

\newrobustcmd{\lorsicper}[2][placeholder]{%
\IfEqCase{#2}{%
{B0p}{\tensor*{\chi}{^{\perp}}}%
{B1p}{\tensor*{\smash{\overset{\wedge}{\chi}}}{^{\perp}_{\indiq[#1]{2}}}}%
{B1m}{\tensor*{\chi}{^{\perp}_{\perp\indiq[#1]{1}}}}%
{B2p}{\tensor*{\smash{\overset{\sim}{\chi}}}{^{\perp}_{\indiq[#1]{2}}}}%
{A0p}{\tensor*{\chi}{^{\perp}_\perp}}%
{A0m}{\tensor*[^{\text{P}}]{\chi}{^{\perp}}}%
{A1p}{\tensor*{\smash{\overset{\wedge}{\chi}}}{^{\perp}_{\perp\indiq[#1]{2}}}}%
{A1m}{\tensor*{\smash{\overset{\rightharpoonup}{\chi}}}{^{\perp}_{\indiq[#1]{1}}}}%
{A2p}{\tensor*{\smash{\overset{\sim}{\chi}}}{^{\perp}_{\perp\indiq[#1]{2}}}}%
{A2m}{\tensor*[^{\text{T}}]{\chi}{^{\perp}_{\indiq[#1]{3}}}}%
}[\packageError{cosmicclass}{Unidentified Critical Case: #1}{}]%
}

\newrobustcmd{\Tl}[2][placeholder]{%
\IfEqCase{#2}{%
{B0p}{\tensor{\chi}{^{\flat}}}%
{B1p}{\tensor{\smash{\overset{\wedge}{\chi}}}{^{\flat}_{\indiq[#1]{2}}}}%
{B1m}{\tensor{\chi}{^{\flat}_{\perp}_{\indiq[#1]{1}}}}%
{B2p}{\tensor{\smash{\overset{\sim}{\chi}}}{^{\flat}_{\indiq[#1]{2}}}}%
{A0p}{\tensor{\chi}{^{\flat}_\perp}}%
{A0m}{\tensor[^{\text{P}}]{\mathcal{T}}{^{\flat}}}%
{A1p}{\tensor{\smash{\overset{\wedge}{\chi}}}{^{\flat}_{\perp\indiq[#1]{2}}}}%
{A1m}{\tensor{\smash{\overset{\rightharpoonup}{\mathcal{T}}}}{^{\flat}_{\indiq[#1]{1}}}}%
{A2p}{\tensor{\smash{\overset{\sim}{\chi}}}{^{\flat}_{\perp\indiq[#1]{2}}}}%
{A2m}{\tensor[^{\text{T}}]{\mathcal{T}}{^{\flat}_{\indiq[#1]{3}}}}%
}[\tensor{\mathcal{T}}{^{\flat}_{\indiq[#1]{3}}}]%
}

\newrobustcmd{\cT}[2][placeholder]{%
\IfEqCase{#2}{%
{B1p}{\tensor{\mathcal{T}}{_{\perp\indiq[#1]{2}}}}%
{B1m}{\tensor{\overset{\rightharpoonup}{\mathcal{T}}}{_{\indiq[#1]{1}}}}%
{A0m}{\tensor[^{\text{P}}]{\mathcal{T}}{}}%
{A2m}{\tensor[^{\text{T}}]{\mathcal{T}}{_{\indiq[#1]{3}}}}%
}[\packageError{cosmicclass}{Unidentified Critical Case: #1}{}]%
}

\newrobustcmd{\cTLambda}[2][placeholder]{%
\IfEqCase{#2}{%
{B1p}{\tensor{\zeta}{_{\perp\indiq[#1]{2}}}}%
{B1m}{\tensor{\overset{\rightharpoonup}{\zeta}}{_{\indiq[#1]{1}}}}%
{A0m}{\tensor[^{\text{P}}]{\zeta}{}}%
{A2m}{\tensor[^{\text{T}}]{\zeta}{_{\indiq[#1]{3}}}}%
}[\packageError{cosmicclass}{Unidentified Critical Case: #1}{}]%
}

\newrobustcmd{\cTpic}[2][placeholder]{%
\IfEqCase{#2}{%
{B1p}{\tensor{\phi}{_{\perp\indiq[#1]{2}}}}%
{B1m}{\tensor{\overset{\rightharpoonup}{\phi}}{_{\indiq[#1]{1}}}}%
{A0m}{\tensor[^{\text{P}}]{\phi}{}}%
{A2m}{\tensor[^{\text{T}}]{\phi}{_{\indiq[#1]{3}}}}%
}[\packageError{cosmicclass}{Unidentified Critical Case: #1}{}]%
}

\newrobustcmd{\cTpicl}[2][placeholder]{%
\IfEqCase{#2}{%
{B1p}{\tensor{\phi}{^{\flat}_{\perp\indiq[#1]{2}}}}%
{B1m}{\tensor{\overset{\rightharpoonup}{\phi}}{^{\flat}_{\indiq[#1]{1}}}}%
{A0m}{\tensor[^{\text{P}}]{\phi}{^{\flat}}}%
{A2m}{\tensor[^{\text{T}}]{\phi}{^{\flat}_{\indiq[#1]{3}}}}%
}[\packageError{cosmicclass}{Unidentified Critical Case: #1}{}]%
}

\newrobustcmd{\cTPiP}[2][placeholder]{%
\IfEqCase{#2}{%
{B1p}{\tensor{\varpi}{_{\perp\indiq[#1]{2}}}}%
{B1m}{\tensor{\overset{\rightharpoonup}{\varpi}}{_{\indiq[#1]{1}}}}%
{A0m}{\tensor[^{\text{P}}]{\varpi}{}}%
{A2m}{\tensor[^{\text{T}}]{\varpi}{_{\indiq[#1]{3}}}}%
}[\packageError{cosmicclass}{Unidentified Critical Case: #1}{}]%
}

\newrobustcmd{\cTl}[2][placeholder]{%
\IfEqCase{#2}{%
{B1p}{\tensor{\mathcal{T}}{^{\flat}_{\perp\indiq[#1]{2}}}}%
{B1m}{\tensor{\smash{\overset{\rightharpoonup}{\mathcal{T}}}}{^{\flat}_{\indiq[#1]{1}}}}%
{A0m}{\tensor[^{\text{P}}]{\mathcal{T}}{^{\flat}}}%
{A2m}{\tensor[^{\text{T}}]{\mathcal{T}}{^{\flat}_{\indiq[#1]{3}}}}%
}[\packageError{cosmicclass}{Unidentified Critical Case: #1}{}]%
}

\newrobustcmd{\cTu}[2][placeholder]{%
\IfEqCase{#2}{%
{B1p}{\tensor{\mathcal{T}}{^{\perp\indiq[#1]{2}}}}%
{B1m}{\tensor{\smash{\overset{\rightharpoonup}{\mathcal{T}}}}{^{\indiq[#1]{1}}}}%
{A0m}{\tensor[^{\text{P}}]{\mathcal{T}}{}}%
{A2m}{\tensor[^{\text{T}}]{\mathcal{T}}{^{\indiq[#1]{3}}}}%
}[\packageError{cosmicclass}{Unidentified Critical Case: #1}{}]%
}

\newrobustcmd{\ncTLambda}[2][placeholder]{%
\IfEqCase{#2}{%
{B0p}{\tensor{\zeta}{^{\indiq[#1]{1}}_{\indiq[#1]{1}\perp}}}%
{B1p}{\tensor{\zeta}{_{[\indiq[#1]{2}]\perp}}}%
{B1m}{\tensor{\zeta}{_{\perp\indiq[#1]{1}\perp}}}%
{B2p}{\tensor{\zeta}{_{\langle\indiq[#1]{2}\rangle\perp}}}%
}[\packageError{cosmicclass}{Unidentified Critical Case: #1}{}]%
}

\newrobustcmd{\ncTmul}[2][placeholder]{%
\IfEqCase{#2}{%
{B0p}{\tensor{\upsilon}{^{\indiq[#1]{1}}_{\indiq[#1]{1}\perp}}}%
{B1p}{\tensor{\upsilon}{_{[\indiq[#1]{2}]\perp}}}%
{B1m}{\tensor{\upsilon}{_{\perp\indiq[#1]{1}\perp}}}%
{B2p}{\tensor{\upsilon}{_{\langle\indiq[#1]{2}\rangle\perp}}}%
}[\packageError{cosmicclass}{Unidentified Critical Case: #1}{}]%
}

\newrobustcmd{\ncTpic}[2][placeholder]{%
\IfEqCase{#2}{%
{B0p}{\tensor{\phi}{^{\indiq[#1]{1}}_{\indiq[#1]{1}\perp}}}%
{B1p}{\tensor{\phi}{_{[\indiq[#1]{2}]\perp}}}%
{B1m}{\tensor{\phi}{_{\perp\indiq[#1]{1}\perp}}}%
{B2p}{\tensor{\phi}{_{\langle\indiq[#1]{2}\rangle\perp}}}%
}[\packageError{cosmicclass}{Unidentified Critical Case: #1}{}]%
}

\newrobustcmd{\ncTpicl}[2][placeholder]{%
\IfEqCase{#2}{%
  {B0p}{\tensor{\phi}{^{\flat}^{\indiq[#1]{1}}_{\indiq[#1]{1}\perp}}}%
{B1p}{\tensor{\phi}{^{\flat}_{[\indiq[#1]{2}]\perp}}}%
{B1m}{\tensor{\phi}{^{\flat}_{\perp\indiq[#1]{1}\perp}}}%
{B2p}{\tensor{\phi}{^{\flat}_{\langle\indiq[#1]{2}\rangle\perp}}}%
}[\packageError{cosmicclass}{Unidentified Critical Case: #1}{}]%
}

\newrobustcmd{\ncTPiP}[2][placeholder]{%
\IfEqCase{#2}{%
{B0p}{\tensor{\varpi}{^{\indiq[#1]{1}}_{\indiq[#1]{1}\perp}}}%
{B1p}{\tensor{\varpi}{_{[\indiq[#1]{2}]\perp}}}%
{B1m}{\tensor{\varpi}{_{\perp\indiq[#1]{1}\perp}}}%
{B2p}{\tensor{\varpi}{_{\langle\indiq[#1]{2}\rangle\perp}}}%
}[\packageError{cosmicclass}{Unidentified Critical Case: #1}{}]%
}

\newrobustcmd{\ncT}[2][placeholder]{%
\IfEqCase{#2}{%
{B0p}{\tensor{\mathcal{T}}{^{\indiq[#1]{1}}_{\indiq[#1]{1}\perp}}}%
{B1p}{\tensor{\mathcal{T}}{_{[\indiq[#1]{2}]\perp}}}%
{B1m}{\tensor{\mathcal{T}}{_{\perp\indiq[#1]{1}\perp}}}%
{B2p}{\tensor{\mathcal{T}}{_{\langle\indiq[#1]{2}\rangle\perp}}}%
}[\packageError{cosmicclass}{Unidentified Critical Case: #1}{}]%
}

\newrobustcmd{\cR}[2][placeholder]{%
\IfEqCase{#2}{%
{A0p}{\tensor{\underline{\mathcal{R}}}{}}%
{A0m}{\tensor[^{\text{P}}]{\mathcal{R}}{_{\perp\circ}}}%
{A1p}{\tensor{\underline{\mathcal{R}}}{_{[\indiq[#1]{2}]}}}%
{A1m}{\tensor{\mathcal{R}}{_{\perp\indiq[#1]{1}}}}%
{A2p}{\tensor{\underline{\mathcal{R}}}{_{\langle\indiq[#1]{2}\rangle}}}%
{A2m}{\tensor[^{\text{T}}]{\mathcal{R}}{_{\perp\indiq[#1]{3}}}}%
}[\packageError{cosmicclass}{Unidentified Critical Case: #1}{}]%
}

\newrobustcmd{\cRLambda}[2][placeholder]{%
\IfEqCase{#2}{%
{A0p}{\tensor{\underline{\zeta}}{}}%
{A0m}{\tensor[^{\text{P}}]{\zeta}{_{\perp\circ}}}%
{A1p}{\tensor{\underline{\zeta}}{_{[\indiq[#1]{2}]}}}%
{A1m}{\tensor{\zeta}{_{\perp\indiq[#1]{1}}}}%
{A2p}{\tensor{\underline{\zeta}}{_{\langle\indiq[#1]{2}\rangle}}}%
{A2m}{\tensor[^{\text{T}}]{\zeta}{_{\perp\indiq[#1]{3}}}}%
}[\packageError{cosmicclass}{Unidentified Critical Case: #1}{}]%
}

\newrobustcmd{\cRl}[2][placeholder]{%
\IfEqCase{#2}{%
  {A0p}{\tensor{\underline{\mathcal{R}}}{^{\flat}}}%
{A0m}{\tensor[^{\text{P}}]{\mathcal{R}}{^{\flat}_{\perp\circ}}}%
{A1p}{\tensor{\underline{\mathcal{R}}}{^{\flat}_{[\indiq[#1]{2}]}}}%
{A1m}{\tensor{\mathcal{R}}{^{\flat}_{\perp\indiq[#1]{1}}}}%
{A2p}{\tensor{\underline{\mathcal{R}}}{^{\flat}_{\langle\indiq[#1]{2}\rangle}}}%
{A2m}{\tensor[^{\text{T}}]{\mathcal{R}}{^{\flat}_{\perp\indiq[#1]{3}}}}%
}[\packageError{cosmicclass}{Unidentified Critical Case: #1}{}]%
}

\newrobustcmd{\cRu}[2][placeholder]{%
\IfEqCase{#2}{%
{A0p}{\tensor{\underline{\mathcal{R}}}{}}%
{A0m}{\tensor[^{\text{P}}]{\mathcal{R}}{_{\perp\circ}}}%
{A1p}{\tensor{\underline{\mathcal{R}}}{^{[\indiq[#1]{2}]}}}%
{A1m}{\tensor{\mathcal{R}}{^{\perp\indiq[#1]{1}}}}%
{A2p}{\tensor{\underline{\mathcal{R}}}{^{\langle\indiq[#1]{2}\rangle}}}%
{A2m}{\tensor[^{\text{T}}]{\mathcal{R}}{^{\perp\indiq[#1]{3}}}}%
}[\packageError{cosmicclass}{Unidentified Critical Case: #1}{}]%
}

\newrobustcmd{\ncR}[2][placeholder]{%
\IfEqCase{#2}{%
  {A0p}{\tensor{\mathcal{R}}{_{\perp\perp}}}%
{A0m}{\tensor[^{\text{P}}]{\mathcal{R}}{_{\circ\perp}}}%
{A1p}{\tensor{\mathcal{R}}{_{\perp[\indiq[#1]{2}]\perp}}}%
{A1m}{\tensor{\mathcal{R}}{_{\indiq[#1]{1}\perp}}}%
{A2p}{\tensor{\mathcal{R}}{_{\perp\langle\indiq[#1]{2}\rangle\perp}}}%
{A2m}{\tensor[^{\text{T}}]{\mathcal{R}}{_{\indiq[#1]{3}\perp}}}%
}[\packageError{cosmicclass}{Unidentified Critical Case: #1}{}]%
}

\newrobustcmd{\ncRLambda}[2][placeholder]{%
\IfEqCase{#2}{%
  {A0p}{\tensor{\zeta}{_{\perp\perp}}}%
{A0m}{\tensor[^{\text{P}}]{\zeta}{_{\circ\perp}}}%
{A1p}{\tensor{\zeta}{_{\perp[\indiq[#1]{2}]\perp}}}%
{A1m}{\tensor{\zeta}{_{\indiq[#1]{1}\perp}}}%
{A2p}{\tensor{\zeta}{_{\perp\langle\indiq[#1]{2}\rangle\perp}}}%
{A2m}{\tensor[^{\text{T}}]{\zeta}{_{\indiq[#1]{3}\perp}}}%
}[\packageError{cosmicclass}{Unidentified Critical Case: #1}{}]%
}

\newrobustcmd{\Proj}[2][placeholder]{%
\IfEqCase{#2}{%
  {A2m}{\tensor[^{\text{T}}]{\check{\mathcal{P}}}{#1}}%
}[\packageError{cosmicclass}{Unidentified Critical Case: #1}{}]%
}

\newrobustcmd{\Projl}[2][placeholder]{%
\IfEqCase{#2}{%
  {A2m}{\tensor[^{\text{T}}]{\check{\mathcal{P}}}{^{\flat}#1}}%
}[\packageError{cosmicclass}{Unidentified Critical Case: #1}{}]%
}

\newcommand{\fcomm}{\varpi}
\newcommand{\icomm}{\vartheta}
\newcommand{\partl}{\mathcal{D}}
%\newcommand{\APiP}{\tensor{\hat{\pi}}{_{ij\overline{k}}}}
\newcommand{\BPiP}{\hat{\pi}}
\newcommand{\nonl}{\hat{\pi}}
\newcommand{\surface}{\partial}


\newrobustcmd{\fA}{%
  {\tensor{\mathcal{  A}}{_{\acu{u}}}}%
}
\newrobustcmd{\fB}{%
  {\tensor{\mathcal{  B}}{_{\acu{v}}}}%
}
\newrobustcmd{\fC}{%
  {\tensor{\mathcal{  C}}{^{\acu{v}}}}%
}
\newrobustcmd{\fphi}{%
  {\tensor{\phi}{^{\acu{w}}}}%
}
\newrobustcmd{\fpi}{%
  {\tensor{\pi}{_{\acu{w}}}}%
}

\newrobustcmd{\covard}[2]{%
  {\frac{\bar{\delta}#1}{\bar{\delta}#2}}
}
\newrobustcmd{\copard}[2]{%
  {\frac{\bar{\partial}#1}{\bar{\partial}#2}}
}
\newrobustcmd{\pard}[2]{%
  {\frac{\partial #1}{\partial #2}}
}


% Journal names 
\def\aj{\rm{AJ}}
\def\actaa{\rm{Acta Astron.}}
\def\araa{\rm{ARA\&A}}
\def\apj{\rm{ApJ}}
\def\apjl{\rm{ApJ}}
\def\apjs{\rm{ApJS}}
\def\ao{\rm{Appl.~Opt.}}
\def\apss{\rm{Ap\&SS}}
\def\aap{\rm{A\&A}}
\def\aapr{\rm{A\&A~Rev.}}
\def\aaps{\rm{A\&AS}}
\def\azh{\rm{AZh}}
\def\baas{\rm{BAAS}}
\def\bac{\rm{Bull. astr. Inst. Czechosl.}}
\def\caa{\rm{Chinese Astron. Astrophys.}}
\def\cjaa{\rm{Chinese J. Astron. Astrophys.}}
\def\icarus{\rm{Icarus}}
\def\jcap{\rm{J. Cosmology Astropart. Phys.}}
\def\jrasc{\rm{JRASC}}
\def\memras{\rm{MmRAS}}
\def\mnras{\rm{MNRAS}}
\def\na{\rm{New A}}
\def\nar{\rm{New A Rev.}}
\def\pra{\rm{Phys.~Rev.~A}}
\def\prb{\rm{Phys.~Rev.~B}}
\def\prc{\rm{Phys.~Rev.~C}}
\def\prd{\rm{Phys.~Rev.~D}}
\def\pre{\rm{Phys.~Rev.~E}}
\def\prl{\rm{Phys.~Rev.~Lett.}}
\def\pasa{\rm{PASA}}
\def\pasp{\rm{PASP}}
\def\pasj{\rm{PASJ}}
\def\rmxaa{\rm{Rev. Mexicana Astron. Astrofis.}}
\def\qjras{\rm{QJRAS}}
\def\skytel{\rm{S\&T}}
\def\solphys{\rm{Sol.~Phys.}}
\def\sovast{\rm{Soviet~Ast.}}
\def\ssr{\rm{Space~Sci.~Rev.}}
\def\zap{\rm{ZAp}}
\def\nat{\rm{Nature}}
\def\iaucirc{\rm{IAU~Circ.}}
\def\aplett{\rm{Astrophys.~Lett.}}
\def\astropj{\rm{Astrophys.~J.}}
\def\apspr{\rm{Astrophys.~Space~Phys.~Res.}}
\def\bain{\rm{Bull.~Astron.~Inst.~Netherlands}}
\def\fcp{\rm{Fund.~Cosmic~Phys.}}
\def\gca{\rm{Geochim.~Cosmochim.~Acta}}
\def\grl{\rm{Geophys.~Res.~Lett.}}
\def\jcp{\rm{J.~Chem.~Phys.}}
\def\jgr{\rm{J.~Geophys.~Res.}}
\def\jqsrt{\rm{J.~Quant.~Spec.~Radiat.~Transf.}}
\def\memsai{\rm{Mem.~Soc.~Astron.~Italiana}}
\def\nphysa{\rm{Nucl.~Phys.~A}}
\def\nphys{\rm{Nucl.~Phys.}}
\def\physrep{\rm{Phys.~Rep.}}
\def\physscr{\rm{Phys.~Scr}}
\def\planss{\rm{Planet.~Space~Sci.}}
\def\procspie{\rm{Proc.~SPIE}}


\def\procrsa{\rm{Proc.~Math.~Phys.~Eng.~Sci.}}
\def\annmath{\rm{Ann.~Math.}}
\def\mathann{\rm{Math.~Ann.}}
\def\annphys{\rm{Ann.~Physics (N.~Y.)}}
\def\annderphys{\rm{Ann.~Phys.}}
\def\ptrsla{\rm{Philos.~Trans.~R.~Soc~A}}
\def\ptrslo{\rm{Philos.~Trans.~R.~Soc~I}}
\def\jmp{\rm{J.~Math.~Phys.}}
\def\cqg{\rm{Class.~Quantum~Gravity}}
\def\physrev{\rm{Phys.~Rev.}}
\def\foundphys{\rm{Found.~Phys.}}
\def\epjc{\rm{Eur.~Phys.~J.~C}}
\def\natast{\rm{Nat.~Astron.}}
\def\pdu{\rm{Phys.~Dark~Universe}}
\def\araa{\rm{Annu.~Rev.~Astron.~Astrophys.}}
\def\nrp{\rm{Nat.~Rev.~Phys.}}
\def\rpp{\rm{Rep.~Prog.~Phys.}}
\def\ijmpd{\rm{Int.~J.~Mod.~Phys.~D}}
\def\pletta{\rm{Phys.~Lett.~A}}
\def\plettb{\rm{Phys.~Lett.~B}}
\def\grg{\rm{Gen.~Relativ.~Gravit.}}
\def\phyrep{\rm{Phys.~Rep.}}
\def\jhep{\rm{J.~High~Energy~Phys.}}
\def\rmphy{\rm{Rev.~Mod.~Phys.~D}}
\def\rmphyo{\rm{Rev.~Mod.~Phys}}
\def\fass{\rm{Front.~Astron.~Space~Sci.}}
\def\progtp{\rm{Prog.~Theor.~Phys.}}
\def\appb{\rm{Acta~Phys.~Pol.~B}}
\def\crphys{\rm{C.~R.~Phys.}}
\def\ijtp{\rm{Int.~J.~Theor.~Phys.}}
\def\progpnp{\rm{Prog.~Part.~Nucl.~Phys.}}
\def\nucphyb{\rm{Nucl.~Phys.~B}}
\def\npcs{\rm{Nonlin.~Phenom.~Complex~Syst.}}
\def\nar{\rm{New.~Astron.~Rev.}}
\def\cjp{\rm{Chin.~J.~Phys.}}
\def\ujp{\rm{Ukr.~J.~Phys.}}
\def\aaca{\rm{Adv.~Appl.~Clifford~Algebras}}
\def\sjetpl{\rm{Sov.~JETP~Lett.}}
\def\mathz{\rm{Math.~Z.}}
\def\macsp{\rm{Mem.~Acad.~St.~Petersbourg}}
\def\rprogphy{\rm{Rept.~Prog.~Phys.}}
\def\annhp{\rm{Ann.~Henri~Poincar\'e}}%Annales de l'I.H.P. Physique th\'eorique
\def\astronj{\rm{Astron.~J.}}
\def\ijgmmp{\rm{Int.~J.~Geom.~Meth.~Mod.~Phys.}}
\def\pispf{\rm{Proc.~Int.~Sch.~Phys.~Fermi}}
\def\livrevrel{\rm{Living~Rev.~Relativ.}}



%Soviet Journal of Experimental and Theoretical Physics Letters

\DeclareRobustCommand{\particle}[1]{%
  \begingroup\normalfont
  %\includegraphics[height=0.4cm]{#1}%
  \includegraphics[height=\fontcharht\font`\B]{#1}%
  \endgroup
}

\DeclareRobustCommand{\ppmfish}{%
  \begingroup\normalfont
  %\includegraphics[height=0.4cm]{#1}%
  \includegraphics[height=0.7\fontcharht\font`\B]{fish_2}%
  \endgroup
}

\newrobustcmd{\PPM}[1]{%
  {\left[\tensor*{\mathsf{M}}{_{\ }^{\left(\text{#1}\right)}}\right]}%
}


\tikzset{
  good/.style={circle, opacity=0.7, draw=green!60, fill=green!5, line width=.8mm, minimum size=3.5mm},
  bad/.style={circle, opacity=0.7, draw=red!60, fill=red!5, line width=.8mm, minimum size=3.5mm},
  badx/.style={circle, opacity=0.4, draw=red!60, fill=red!5, line width=.8mm, minimum size=3.5mm},
  save/.style={circle, opacity=0.7, draw=green!60, fill=green!60, line width=1.5mm, minimum size=3.5mm},
  kill/.style={circle, opacity=0.7, draw=red!60, fill=red!60, line width=1.5mm, minimum size=3.5mm},
  bind/.style={draw=green!60, opacity=0.7, line width=1mm,},
  wrap/.style={draw=red!60, opacity=0.7, line width=1mm,},
}

 %some personal macros

\usepackage{hyperref}
\hypersetup{%
     colorlinks = true,%
     linkcolor = Blue,%
     citecolor = Blue,%
     filecolor = Blue,%
     urlcolor = Blue% 
     }%
\usepackage[capitalize]{cleveref} %always load this last in preamble

\begin{document}

\title{A Purely Gravitational Origin for Einstein-Proca Theory}

\author{Will Barker}
\email{wb263@cam.ac.uk}
\affiliation{Astrophysics Group, Cavendish Laboratory, JJ Thomson Avenue, Cambridge CB3 0HE, UK}
\affiliation{Kavli Institute for Cosmology, Madingley Road, Cambridge CB3 0HA, UK}
\author{Sebastian Zell}
\email{sebastian.zell@uclouvain.be}
\affiliation{Centre for Cosmology, Particle Physics and Phenomenology -- CP3,
	Universit\'e catholique de Louvain, B-1348 Louvain-la-Neuve, Belgium}

%\date{}

\begin{abstract}
	We construct a theory of gravity in which a propagating massive vector field arises from a quadratic curvature invariant. 
	The Einstein-Cartan formulation and a partial suppression of torsion ensure the absence of ghost and strong-coupling problems, as we prove with nonlinear Lagrangian and Hamiltonian analysis. Augmenting General Relativity with a propagating torsion vector, our theory provides a purely gravitational origin of Einstein-Proca models and constrains their parameter space. As an outlook to phenomenology, we discuss the gravitational production of fermionic dark matter.
\end{abstract}



\maketitle


\section{Open questions in General Relativity}
The theory of General Relativity (GR) provides an outstandingly successful description of gravity and has been confirmed by countless experiments, such as the breakthrough discovery of gravitational waves~\cite{LIGOScientific:2016aoc}. Nevertheless, many far-reaching questions have remained unanswered. Gravitational interactions have long revealed the existence of a non-baryonic form of matter~\cite{Zwicky:1933gu}. It is unknown, however, what this dark matter is composed of and how it was produced in the cosmological history~\cite{Feng:2010gw,Peebles:2013hla,Bertone:2016nfn,Arbey:2021gdg}. Moreover, the very early moments of our Universe -- due to a suspected phase of inflation~\cite{Starobinsky:1980te, Guth:1980zm, Linde:1981mu, Mukhanov:1981xt} -- and its far future -- because of the observed dark energy~\cite{SupernovaSearchTeam:1998fmf, SupernovaCosmologyProject:1998vns} -- appear to have in common an accelerated expansion. Agreement with recent	measurements, in particular of the cosmic microwave background (CMB)~\cite{Planck:2018jri, BICEP:2021xfz}, is excellent but a microscopic origin has still not been determined both for dark energy and for inflation.

It is exciting to explore if GR also contains an answer to these problems.  As we shall show, this question is closely connected to the geometry of spacetime, which a priori is characterized by three independent properties curvature $\rcR{}$, torsion $T^\alpha_{~~\beta \gamma}$ and nonmetricity $Q_{\alpha\mu\nu}$ (see~\cref{fig:schema}). The presence or absence of $\rcR{}$, $Q_{\alpha\mu\nu}$ and $T^\alpha_{~~\beta \gamma}$ defines seven formulations of GR~\cite{Einstein:1915,Weyl:1918, Palatini:1919,Weyl:1922, Eddington:1923, Cartan:1922, Cartan:1923,Eddington:1923, Cartan:1924, Cartan:1925,Einstein:1925, Einstein:1928, Einstein:19282,Schroedinger:1950,Moller:1961, Pellegrini:1963, Hayashi:1967se,Cho:1975dh,Hehl:1976kt, Hehl:1976kv, Hehl:1976my, Hehl:1977fj,Kijowski:1978,Hayashi:1979qx, Nester:1998mp, BeltranJimenez:2019odq} (see~\cite{Heisenberg:2018vsk, BeltranJimenez:2019esp, Rigouzzo:2022yan} for an overview\footnote{Additionally, there is purely affine gravity~\cite{Eddington:1923,Einstein:1925,Schroedinger:1950, Kijowski:1978}, where $\Gamma^\alpha_{~~\beta \gamma}$ is the only dynamical field.}). For example, assuming that both $Q_{\alpha\mu\nu}$ and $T^\alpha_{~~\beta \gamma}$ vanish leads to the commonly-used metric variant of GR~\cite{Einstein:1915} corresponding to Riemannian geometry, whereas Einstein-Cartan (EC) gravity is defined by including both $\rcR{}$ and $T^\alpha_{~~\beta \gamma}$ while excluding $Q_{\alpha\mu\nu}$~\cite{Cartan:1922,Cartan:1923,Cartan:1924,Cartan:1925,Einstein:1925, Einstein:1928,Einstein:19282}. The EC version stands out since it can be derived by gauging the Poincar\'e group~\cite{Utiyama:1956sy, Sciama:1962, Kibble:1961ba}, which puts gravity on the same footing as the other fundamental forces. At first sight, the various versions of GR appear to be very different. In the classical, matter-free limit, however, they are all fully equivalent.

	\begin{figure}
		\vspace{-0.5\baselineskip}
	\includegraphics[width=1\linewidth]{NonRiemannianSchematic.pdf}
		\vspace{-2.5\baselineskip}
		\caption{Schematic representation of parallel transport in the presence of curvature $\tensor{R}{^\rho_{\sigma\mu\nu}}$ -- rotation around closed curves -- torsion $\tensor{T}{^\sigma_{\mu\nu}}$ -- non-closure of (infinitesimal) parallelograms -- and nonmetricity $\tensor{Q}{_{\lambda\mu\nu}}$ -- non-conservation of vector norms.}
	\label{fig:schema}
\end{figure}

There are two main ways to break this equivalence. First, one can couple matter (non-minimally) to GR. In this case, the spectrum of gravity remains unchanged, \ie no additional propagating \dof emerge apart from the massless spin-2 graviton, but new effects arise at high energies~\cite{Freidel:2005sn,Bauer:2008zj,Poplawski:2011xf,Diakonov:2011fs,Khriplovich:2012xg,Magueijo:2012ug,Khriplovich:2013tqa,Markkanen:2017tun,Carrilho:2018ffi,Enckell:2018hmo,Rasanen:2018fom,Rubio:2019ypq,Shaposhnikov:2020geh,Karananas:2020qkp,Langvik:2020nrs,Shaposhnikov:2020gts,Mikura:2020qhc,Shaposhnikov:2020aen,Kubota:2020ehu,Enckell:2020lvn,Iosifidis:2021iuw,Racioppi:2021ynx,Cheong:2021kyc,Dioguardi:2021fmr,Piani:2022gon,Dux:2022kuk,Rasanen:2022ijc,Gialamas:2023emn,Gialamas:2023flv,Piani:2023aof,Poisson:2023tja,Rigouzzo:2023sbb}. An example, which we shall discuss later, is as a novel mechanism for producing dark matter~\cite{Shaposhnikov:2020aen}.  Second, one can add terms nonlinear in curvature. Almost inevitably, this leads to the existence of new dynamical particles~\cite{Stelle:1977ry,Neville:1978bk,Neville:1979rb,Sezgin:1979zf,Hayashi:1979wj, Hayashi:1980qp}. As discussed in more detail below, these additional propagating \dof are generically plagued by various types of inconsistencies~\cite{Stelle:1977ry,Neville:1978bk,Neville:1979rb,Sezgin:1979zf,Hayashi:1979wj, Hayashi:1980qp,Hecht:1990wn,Hecht:1991jh,Chen:1998ad}. Finding a model which is consistent and weakly coupled has turned out to be challenging and, for torsion, so far was only
achieved for a propagating scalar mode~\cite{Hecht:1996np,Chen:1998ad,Yo:1999ex,Yo:2001sy}.



In this paper, we will for the first time obtain a propagating torsion vector from a curvature-squared term in a model for which consistency can be proven. To construct our theory, we shall combine the EC formulation of GR with the concept of partial Lagrange multipliers~\cite{Barker:2022jsh}. In addition to a proof of concept for a new class of models, our findings have important implications for phenomenology since dynamical (Abelian) massive vector fields have been employed in various contexts, in particular for inflation~\cite{Ford:1989me,Golovnev:2008cf,Chiba:2008eh,Koivisto:2008xf,Kanno:2008gn,Emami:2016ldl}, dark energy~\cite{Armendariz-Picon:2004say,Koivisto:2007bp,Koivisto:2008xf,Tasinato:2014eka,DeFelice:2016yws,DeFelice:2016uil,BeltranJimenez:2016afo,deFelice:2017paw,Nakamura:2018oyy,Heisenberg:2020xak,Benisty:2021sul,deRham:2021efp} and dark matter~\cite{Hambye:2008bq,Hambye:2009fg,Arina:2009uq,Hisano:2010yh,Diaz-Cruz:2010czr,Lebedev:2011iq,Farzan:2012hh,Baek:2012se,Arcadi:2020jqf,Barman:2021qds} (see also~\cite{Bekenstein:1971hc,Heisenberg:2017hwb,Garcia-Saenz:2021uyv} for implications for black holes). In these works, however, an a priori independent vector field was considered as an ad hoc addition to GR. Our theory endows these Einstein-Proca models with a gravitational origin. This is not only conceptually appealing but also leads to constraints on parameters, which would otherwise seem arbitrary from an effective field theory perspective. Moreover, we show that the massive vector mode can be integrated out in relevant parts of parameter space, thereby establishing a connection to works with non-propagating torsion using the example of~\cite{Shaposhnikov:2020aen}.

\section{Obstructions to dynamical torsion}
The simplest choice of action for EC gravity is (see~\cite{Hehl:1976kj, Shapiro:2001rz, Blagojevic:2002, Blagojevic:2012bc} for reviews, and all our conventions in~\cref{Conventions})
\begin{equation} \label{EC}
	L_{\text{EC}}  \equiv -\frac{1}{2} \Planck^2 \rcR{} 	+ \TLambda[Q]{_{\mu\nu\sigma}}Q{^{\mu\nu\sigma}}+L_{\text{M}}\left(\Gamma\right) \;,
\end{equation}
where the multiplier field $\TLambda[Q]{_{\mu\nu\sigma}}$ enforces the vanishing of nonmetricity and we included matter $L_{\text{M}}\left(\Gamma\right)$. We can decompose curvature in its Riemannian part $\rR{}$ and contributions from irreducible parts of torsion:
	\begin{align}
		\rcR{}\left(Q_{\alpha\mu\nu}\to 0\right) & = \rR{}
		+
		\frac{8}{9}
		\T[1]{_{\mu\nu\sigma}}
		\T[1]{^{\mu[\nu\sigma]}}
		-
		\frac{2}{3}
		\T[2]{_{\mu}}
		\T[2]{^{\mu}}	
		\nonumber\\
		&\ \ \ 
		+
		\frac{3}{2}
		\T[3]{_{\mu}}
		\T[3]{^{\mu}}	
		-
		2\rD{_\mu}\T[2]{^\mu} \;.
\label{decomposition}
	\end{align}
Varying~\cref{EC} \wrt torsion and without matter, it follows that $\T[1]{^{\mu\nu\sigma}}\approx\T[2]{^{\mu}}\approx\T[3]{^{\mu}}\approx 0$. Plugging this result back in the action shows that the theory \eqref{EC} is equivalent to the metric formulation of GR, $L_{\text{GR}} \equiv -\frac{1}{2} \Planck^2 \rR{}$. 

Adding terms that are quadratic in curvature~\cite{Utiyama:1956sy,Sciama:1964wt,Hayashi:1980ir,Hayashi:1980qp} to~\eqref{EC} generically leads to two types of problems caused by additional propagating d.o.f. The first one consists of ghost and tachyonic instabilities, \ie kinetic or mass terms with the wrong sign~\cite{Stelle:1977ry,Neville:1978bk,Neville:1979rb,Sezgin:1979zf,Hayashi:1979wj, Hayashi:1980qp}. Superficially, one can eliminate these inconsistencies by tuning the kernel of the linearised wave operator, and indeed a `zoo' of such unitary cases is available~\cite{Sezgin:1981xs,Kuhfuss:1986rb,Yo:1999ex,Yo:2001sy,Nair:2008yh,Nikiforova:2009qr,Karananas:2014pxa,Karananas:2016ltn,Obukhov:2017pxa,Blagojevic:2017ssv,Blagojevic:2018dpz,Lin:2018awc,Jimenez:2019qjc,Lin:2019ugq,BeltranJimenez:2019acz,Aoki:2019rvi,Percacci:2019hxn,BeltranJimenez:2020sqf,Marzo:2021esg,Marzo:2021iok,Baldazzi:2021kaf,Blagojevic:2002,Blagojevic:1983zz,Blagojevic:1986dm,Blagojevic:2018dpz,Barker:2021oez,Ong:2013qja,Blagojevic:2018dpz,Karananas:2014pxa,delaCruzDombriz:2021nrg,MaldonadoTorralba:2020mbh,Zhang:2019xek,Zhang:2019mhd,Tseng:2018feo,Ho:2015ulu,Yo:2006qs,Shie:2008ms,Chen:2009at,Baekler:2010fr,Ho:2011qn,Ho:2011xf,Puetzfeld:2004yg,Puetzfeld:2014sja,Barker:2020gcp,Ni:2009fg,Ni:2015poa,Annala:2022gtl}. However, a second class of issues, first discussed in the metric-compatible teleparallel formulation of GR~\cite{Einstein:1928, Einstein:19282, Moller:1961, Pellegrini:1963, Hayashi:1967se,Cho:1975dh, Hayashi:1979qx} and its generalization in `\emph{New GR}' (NGR)~\cite{Hayashi:1979qx} (see~\cite{Hecht:1991jh,Magueijo:2008sx,Afshordi:2006ad,Dimakis:1989az,Dimakis:1989ba,Lemke:1990su,Hecht:1990wn}), is only visible in a nonlinear analysis~\cite{Delhom:2022vae,Velo:1969txo,Aragone:1971kh,Krasnov:2021zen,Chen:1998ad,Cheng:1988zg,Chen:1998ad,Blixt:2018znp,Blixt:2019ene,Blixt:2020ekl}. Applied to EC gravity, an especially restrictive report~\cite{Yo:2001sy} found that for propagating tensor/vector torsion, weak-field linearisation causes over-counting of the constraints~\cite{Blagojevic:2002,Blagojevic:1983zz,Blagojevic:1986dm,Blagojevic:2018dpz,Barker:2021oez,Ong:2013qja} -- peculiar to torsion theories, the activated fields are ghosts whenever the linear spectrum is unitary~\cite{Yo:2001sy,Hayashi:1980qp,Blagojevic:1983zz,Barker:2021oez,Blagojevic:2013dea,Blagojevic:2013taa,Aoki:2020rae,BeltranJimenez:2019hrm}. In the modern parlance (see~\cref{SCIntro},~\cite{Delhom:2022vae,Baumann:2011dt,Barker:2022kdk,JimenezCano:2021rlu,Annala:2022gtl,Yo:2001sy,Baumann:2011dt,Charmousis:2008ce,Wang:2017brl,Charmousis:2009tc,Papazoglou:2009fj}), this means that these new particles become \emph{strongly coupled} on Minkowski spacetime, which is thereby rendered dynamically unreachable~\cite{BeltranJimenez:2020lee,Barker:2022kdk,Mazuet:2017rgq,Gumrukcuoglu:2012aa,DAmico:2011eto}.\footnote{Such issues of strong coupling have been discovered in a variety of theories~\cite{Charmousis:2008ce,Wang:2017brl,Charmousis:2009tc,Papazoglou:2009fj,Vainshtein:1972sx,Deffayet:2001uk,Deffayet:2005ys,deRham:2014zqa,Deser:2014hga}.} As stated above, models that avoid this problem and remain weakly coupled have so far only been constructed for a propagating \emph{scalar} mode of torsion~\cite{Hecht:1996np,Yo:1999ex,Yo:2001sy} (see e.g.~\cite{delaCruzDombriz:2021nrg,MaldonadoTorralba:2020mbh,Zhang:2019xek,Zhang:2019mhd,Tseng:2018feo,Ho:2015ulu,Yo:2006qs,Shie:2008ms,Chen:2009at,Baekler:2010fr,Ho:2011qn,Ho:2011xf} and~\cite{Puetzfeld:2004yg,Puetzfeld:2014sja,Barker:2020gcp,Ni:2009fg,Ni:2015poa} for applications and reviews, respectively, and~\cite{Jimenez-Cano:2022sds,Iosifidis:2023pvz,Marzo:2021esg,Marzo:2021iok,Piva:2021nyj,Iosifidis:2021xdx,Percacci:2020ddy,BeltranJimenez:2020sqf,BeltranJimenez:2019acz,BeltranJimenez:2019esp} for similar prospects in the nonmetric sector).

In this paper we construct a healthy theory with \emph{propagating vector} (PV) torsion. Extending~\cref{EC} it is defined\footnote{We can also now put forward \emph{vector torsion} on an equal footing with Weyl's vector nonmetricity~\cite{Iosifidis:2018zwo,BeltranJimenez:2014iie,Iosifidis:2018diy,Helpin:2019vrv,OrejuelaGarcia:2020viw,Ghilencea:2020piz,BeltranJimenez:2020sih,Xu:2020yeg,Yang:2021fjy,Quiros:2021eju,Quiros:2022uns,Yang:2022icz,Burikham:2023bil,Haghani:2023nrm} -- indeed our $\rcR{_{[\mu\nu]}}$ in~\cref{DC} is analogous to Weyl's homothetic curvature~\cite{BeltranJimenez:2016wxw,Iosifidis:2018jwu}.}
\begin{equation}\label{DC}
\begin{aligned}
\LPV{}&\equiv-\frac{1}{2}\Planck^2\rcR{}
+\alp{5}\rcR{_{[\mu\nu]}}
\rcR{^{[\mu\nu]}}
+\bet{2}\Planck^2\T[2]{_{\mu}}\T[2]{^{\mu}}
\\  
	&\hspace{-25pt} 
+\TLambda[1]{_{\mu}^{\nu\sigma}}\T[1]{^{\mu}_{\nu\sigma}}
+\TLambda[3]{^\mu}\T[3]{_\mu}
+\TLambda[Q]{_{\mu\nu\sigma}}Q{^{\mu\nu\sigma}}
+L_{\text{M}}\left(\Gamma\right) \;,
\end{aligned}
\end{equation}
where $\TLambda[1]{_{\mu}^{\nu\sigma}}$ and $\TLambda[3]{^\mu}$ are two Lagrange multipliers. Instead of postulating the existence of a kinetic term for $\T[2]{_\mu}$ (see \eg~\cite{Carroll:1994dq, Shapiro:2001rz, Barman:2019mlj} and also~\cite{Vollick:2006uq}), we shall show that it arises from $\rcR{_{[\mu\nu]}} \rcR{^{[\mu\nu]}}$.
Unlike in recent attempts~\cite{Jimenez-Cano:2022sds,Iosifidis:2023pvz,BeltranJimenez:2020sqf,BeltranJimenez:2019acz,BeltranJimenez:2019esp}, we will demonstrate how consistent dynamics and unitarity are simultaneously upheld by \emph{complete} nonlinear analysis. 

\section{Singular spectrum without multipliers} 
Our key innovation in our~\cref{DC} is the \emph{partial} use of multiplier fields~\cite{Barker:2022jsh}. To motivate them, consider what happens if they are omitted. For simplicity we take $L=\LPV{\TLambda[1]{_{\mu}^{\nu\sigma}},\TLambda[3]{^{\mu}},\bet{2},{L_{\mathrm{M}}}\to 0}$.\footnote{It can be shown that for non-zero $\bet{2}$ only one of the two torsion vectors is strongly coupled.} We package six \dof into the 2-form field 
\begin{equation}\label{KalbRamond}
	\begin{aligned}
		\B{_{\mu\nu}}&\equiv 
		\Planck^{-1}\Big(2\rD{_{\sigma}}\T[1]{^{\sigma}_{[\mu\nu]}}
	-\F[2]{_{\mu\nu}}
	+\frac{3}{4}\tensor{\epsilon}{^{\sigma\lambda}_{\mu\nu}}\F[3]{_{\sigma\lambda}}\\
		&\hspace{20pt}
		+2\T[2]{_{\sigma}}\T[1]{^{\sigma}_{[\mu\nu]}}
		+3\tensor{\epsilon}{^{\sigma\lambda\rho}_{[\mu}}\T[3]{_{\sigma}}\T[1]{_{\nu]\lambda\rho}}\Big) \;,
	\end{aligned}
\end{equation}
where the Maxwell field strengths are $\F[2]{_{\mu\nu}}\equiv 2\rD{_{[\mu}}\T[2]{_{\nu]}}$ etc. We introduce~\cref{KalbRamond} just to simplify the $\tensor{\omega}{^{ij}_\mu}$-equations $\delta/\delta \tensor{\omega}{^{ij}_\mu}\int\mathrm{d}^4xeL\approx 0$, which decompose to~\cite{FieldEquations}
\begin{subequations}
	\begin{align}
		%&\hspace{-30pt}
		\T[1]{^\sigma_{\mu\nu}}
		&\approx
		\frac{\alp{5}}{\Planck}\left[
		\OmitIndices{\rD{}\B{}}
		+\frac{\OmitIndices{\B{}\rD{}\B{}}}{\Planck}
		+\frac{\OmitIndices{\B{}^2\T[1]{}}}{\Planck}
		\right] \;,\label{TensorEquation}\\
		\T[2]{_\mu}
		&\approx
		\frac{4\alp{5}}{3\Planck}\Big[\rD{_\nu}\B{_\mu^\nu}
		-\B{_{\sigma\lambda}}\T[1]{_{\mu}^{\sigma\lambda}}\Big]\;,\label{VectorEquation}\\
		\T[3]{_\mu}
		&\approx
		\frac{4\alp{5}\tensor{\epsilon}{_{\mu}^{\nu\sigma\lambda}}}{9\Planck}\Big[\rD{_\nu}\B{_{\sigma\lambda}}
		-2\B{^{\rho}_{\nu}}\T[1]{_{\rho\sigma\lambda}}\Big]\;,\label{AxialVectorEquation}
	\end{align}
\end{subequations}
where $\OmitIndices{...}$ suppresses contractions, but all parts are simplified by~\eqref{KalbRamond}. With $\alp{5}\to 0$ we recover entirely vanishing vacuum torsion as expected in EC theory, otherwise~\crefrange{TensorEquation}{AxialVectorEquation} should be wavelike for dynamical torsion (if any). It is simplest to notice how all torsion dynamics can be confined to $\B{_{\mu\nu}}$, though that variable eliminates a single derivative in~\cref{KalbRamond}. To extract the propagating (second-derivative) equation in $\B{_{\mu\nu}}$, we take the antisymmetrised divergence of~\eqref{TensorEquation}, next eliminating $\rD{_{\sigma}}\T[1]{^{\sigma}_{[\mu\nu]}}$ for $\B{_{\mu\nu}}$, $\T[1]{^\mu_{\nu\sigma}}$, $\T[2]{_\mu}$ and $\T[3]{_\mu}$ using~\cref{KalbRamond}, then using~\cref{VectorEquation,AxialVectorEquation} to eliminate $\T[2]{_\mu}$ and $\T[3]{_\mu}$ for $\B{_{\mu\nu}}$ and $\T[1]{^\mu_{\nu\sigma}}$, before finally recycling~\cref{TensorEquation} to eliminate all remaining $\T[1]{^\mu_{\nu\sigma}}$ \emph{perturbatively} in terms of $\B{_{\mu\nu}}$~\cite{FieldEquations}. Upon integrating, we find (at least on flat space) that the resulting equation descends from the effective theory 
\begin{align}\label{KalbRamondEffectiveTheory}
	&\LPV{\TLambda[1]{_{\mu}^{\nu\sigma}},\TLambda[3]{^{\mu}},\bet{2},{L_{\mathrm{M}}},\rR{}\to 0}
\nonumber\\
	&\hspace{20pt} \cong 	
	-\frac{\Planck^2}{2}\B{_{\mu\nu}}\B{^{\mu\nu}}
	+\frac{\alp{5}^3}{\Planck^2}\OmitIndices{\B{}^2\rD{}\B{}\rD{}\B{}}
+\mathcal{O}\left(B^6\right) \;.
\end{align}
A pathological kinetic term is revealed in~\cref{KalbRamondEffectiveTheory}, not e.g. the safe $p$-form operator $\rD{_{[\mu}}\tensor{B}{_{\nu\sigma]}}\rD{^{[\mu}}\tensor{B}{^{\nu\sigma]}}$~\cite{Henneaux:1986ht}. To see why, truncate~\eqref{KalbRamondEffectiveTheory} at~$\mathcal{O}\left(B^4\right)$ as shown, then heuristically the \emph{canonical} 2-form has a mass~$\sim\Planck/B$ as~$B\to 0$, whereupon it drops out of the linear spectrum (see~\cref{SCIntro}). Couplings of $B^n$ vertices would similarly scale as $B^{-n}$, and hence become \emph{strong}. The problem is only aggravated at higher perturbative orders. We can compare~\eqref{KalbRamondEffectiveTheory} with the seemingly unrelated NGR model~\cite{Hayashi:1979qx,Moller:1961,Pellegrini:1963}. That theory also has a dynamical 2-form $\B{_{\mu\nu}}\equiv\Planck\tensor{\eta}{_{ij}}\tensor*{\delta}{^i_{[\mu}}\tensor{e}{^j_{\nu]}}$, which is different to~\eqref{KalbRamond}. Instead of diverging as~$B\to 0$, the NGR 2-form mass \emph{vanishes}, but again this \emph{removes} longitudinal polarisations from the linear spectrum~\cite{BeltranJimenez:2019nns,Koivisto:2018loq,BeltranJimenez:2020fvy}. Relative to our~\cref{KalbRamondEffectiveTheory}, NGR is positioned at the far side of the non-Riemannian landscape: yet both theories fail as minimal deviations from the dynamically equivalent EC or teleparallel models respectively.

\section{Healthy spectrum with multipliers} 
Obviously~\cref{KalbRamondEffectiveTheory} is sick: we will now show that our $\TLambda[1]{_{\mu\nu\sigma}}\T[1]{^{\mu\nu\sigma}}$ term is the cure. For greater generality we restore the whole axial vector sector by considering $L=\LPV{\TLambda[3]{^\mu}\to 0}+\bet{3}\Planck^2\T[3]{^\mu}\T[3]{_\mu}$. Using the spin tensor of matter $e\tensor{S}{^\mu_{ij}}\equiv -2\delta/\delta \tensor{\omega}{^{ij}_\mu}\int\mathrm{d}^4xeL_{\text{M}}$,~\crefrange{TensorEquation}{AxialVectorEquation} become
\begin{subequations}
	\begin{align}
		&\hspace{-30pt}
		\Si[1]{^\mu_{\nu\sigma}}
		\approx
		\TLambda[1]{^\mu_{\nu\sigma}}
		+\alp{5}\Big[
			\OmitIndices{\tensor[^{(2,3)}]{T}{}\rD{}\tensor[^{(2,3)}]{T}{}}
		+\OmitIndices{\rR{}\tensor[^{(2,3)}]{T}{}}
		+\dots
		\Big] \;,\label{NewTensorEquation}\\
		\Si[2]{_\mu}
		&\approx
		2\left(1+3\bet{2}\right)\Planck^2\T[2]{_\mu}
		+\frac{8\alp{5}}{3}\rD{_\nu}\F[2]{_\mu^\nu} \;,
		\label{NewVectorEquation}\\
		\Si[3]{_\mu}
		&\approx
		-\left(1-4\bet{3}/3\right)\Planck^2\T[3]{_\mu}
		-\frac{4\alp{5}}{3}\rD{_\nu}\F[3]{_\mu^\nu} \;.
		\label{NewAxialVectorEquation}
	\end{align}
\end{subequations}
But~\cref{NewVectorEquation,NewAxialVectorEquation} are Proca equations, whilst~\cref{NewTensorEquation} eliminates $\TLambda[1]{_\mu^{\nu\sigma}}$ in the (asymmetric) $\tensor{e}{^i_{\mu}}$-equation. The antisymmetric part is then an identity; the symmetric part is the $\tensor{g}{_{\mu\nu}}$-equation of the following effective torsion-free theory, to be compared with~\cref{KalbRamondEffectiveTheory}
\begin{align}\label{MasterLagrangian}
	&\hspace{-15pt}\LPV{\TLambda[3]{^\mu}\to 0}+\bet{3}\Planck^2\T[3]{^\mu}\T[3]{_\mu}
	\nonumber\\
	&
	\cong 
	-\frac{\Planck^2\rR{}}{2}
	+
	\frac{2\alp{5}}{9}
	\F[2]{_{\mu\nu}}
	\F[2]{^{\mu\nu}}
	-
	\frac{\alp{5}}{2}
	\F[3]{_{\mu\nu}}
	\F[3]{^{\mu\nu}}
	\nonumber\\
	&
	+\Planck^2
	\left(
	\tfrac{\left(1+3\bet{2}\right)}{3}
	\T[2]{_{\mu}}
	\T[2]{^{\mu}}	
	-
	\tfrac{\left(3-4\bet{3}\right)}{4}
	\T[3]{_{\mu}}
	\T[3]{^{\mu}}	
	\right)
	\nonumber\\
	&
	-
	\frac{1}{3}
	\T[2]{_{\mu}}
	\Si[2]{^{\mu}}	
	-
	\frac{3}{2}
	\T[3]{_{\mu}}
	\Si[3]{^{\mu}}
	+L_{\text{M}}\left(\rCon{}\right) \;.
\end{align}
In~\cref{MasterLagrangian}, we notice for full consistency that the residual torsion reduces to the Proca pair in~\cref{NewVectorEquation,NewAxialVectorEquation}, one of which is a ghost, and the full model $\LPV{}$ or $\LPV{(2)\rightleftarrows(3)}$ kills off the ghost in either case. In contrast, it is critical to understand that $\TLambda[1]{^\lambda_{\mu\nu}}$ does not merely \emph{kill off} the strongly coupled modes: both~\cref{KalbRamondEffectiveTheory} and~\cref{MasterLagrangian} propagate six extra non-graviton \dof, so $\TLambda[1]{^\lambda_{\mu\nu}}$ weakens the strong coupling.
In~\cref{DC}, valid for $\alp{5}<0$, the mass of $\T[2]{_\mu}$ is 
\begin{equation}
	\mass[2]{}^2\equiv-3\Planck^2(1+3\bet{2})/4\alp{5} \;.\label{Mass2}
\end{equation}
Perfectly analogous results hold for $\alp{5}>0$, and the $\T[3]{_\mu}$ mass is $\mass[3]{}^2\equiv-\Planck^2\left(3-4\bet{3}\right)/4\alp{5}$.

We shall briefly comment on how our model \eqref{DC} relates to the known formulations of GR. Since some of the irreducible representations of torsion vanish because of partial multipliers~\cite{Barker:2022jsh}, one could classify it as interpolating between the metric and EC formulations, as illustrated in~\cref{Propaganda}. However, our theory \eqref{DC} also shares properties with teleparallel versions of GR, in which curvature is excluded~\cite{Einstein:1928, Einstein:19282, Moller:1961, Pellegrini:1963, Hayashi:1967se,Cho:1975dh, Hayashi:1979qx, Nester:1998mp, BeltranJimenez:2019odq}. The reason is that -- as shown -- the multiplier $\TLambda[1]{_{\mu}^{\nu\sigma}}$ changes the spectrum. The same happens in teleparallel theories for the Lagrange multiplier that enforces the vanishing of curvature (see~\cite{BeltranJimenez:2018vdo,Rigouzzo:2022yan}).

\begin{figure}[t]
  \center
  \includegraphics[width=\linewidth]{PropagandaLite.pdf}
	\caption{\label{Propaganda} 
	Multipliers are already ubiquitous in gravity. Placement of our PV theory~\cref{DC}, relative to Einstein-Cartan (EC)~\cite{Cartan:1922,Cartan:1923,Cartan:1924,Cartan:1925,Einstein:1925, Einstein:1928,Einstein:19282} and metric GR~\cite{Einstein:1915}.
	}
\end{figure}

\section{Illustration of strong coupling alleviation} 
We shall illustrate the mechanism by which $\TLambda[1]{_{\mu}^{\nu\sigma}}$ alleviates the strong coupling problem in a simple example -- the corresponding analysis for our theory \eqref{DC} is presented in~\cref{FullAnalysis}. Assuming familiarity with the Dirac algorithm~\cite{Anderson:1951ta,Bergmann:1954tc,Castellani:1981us,Yo:2001sy,1999IJMPD...8..459Y,Chen:1998ad} (see good pedagogical introductions~\cite{Golovnev:2022rui,Henneaux:1992ig,Blagojevic:2002}), we consider analogs of `\emph{tetrad}' $\qe{}$, `\emph{spin connection}' $\qw{}$ and `\emph{torsion}' $\dqe{}+\qw{}$ in a minimal working example (MWE)
\begin{equation}\label{MWE}
	\begin{aligned}
		L_{\text{MWE}}&\equiv\qw{}\left[\qe{}\left(1+\qw{}\right)-1\right]\left(\dqw{}+1\right)\\
			&\ \ \ \ \ +\left(\qe{}-1\right)^2+\lambda\left(\dqe{}+\qw{}\right) \;,
	\end{aligned}
\end{equation}
with `\emph{Minkowski background}' $\qe{}-1\approx\qw{}\approx0$.
First try omitting the multiplier: the definitions $p_i\equiv\partial/\partial\dot{q}_i L_{\text{MWE}}\left(\lambda\to 0\right)$ engender two \emph{primary} constraints $\phie{}\equiv\pie{}\approx 0$ and $\phiw{}\equiv\piw{}-\qw{}\left[\qe{}\left(1+\qw{}\right)-1\right]\approx 0$, so the Hamiltonian is $H_{\text{MWE}}\left(\lambda\to 0\right)\equiv\sum_i u_i\phi_i-\qw{}\left[\qe{}\left(1+\qw{}\right)-1\right]-\left(\qe{}-1\right)^2$, where $u_i$ replace Dirac's missing/uninvertible velocities $\dot{q}_i$. The Poisson bracket $\left\{\phie{},\phiw{}\right\}\approx \qw{}\left(1+\qw{}\right)$ has the sickly property that it \emph{vanishes} in the background linearisation of~\cref{MWE}. This feature is the popular understanding of strongly coupled torsion~\cite{Yo:2001sy,Barker:2021oez,Cheng:1988zg,Chen:1998ad,Blixt:2020ekl,1999IJMPD...8..459Y}. For a vanishing bracket, the consistency conditions $\dot{\phi}_i\equiv\left\{\phi_i,H_{\text{MWE}}\right\}\approx0$ engender \emph{secondary} constraints $\chie{}\equiv 2\left(\qe{}-1\right)+\qw{}\left(\qw{}+1\right)\approx 0$ and $\chiw{}\equiv\qe{}\left(1+2\qw{}\right)-1\approx 0$. Since $\{\chie{},\phie{}\}\approx 2$ and $\{\chiw{},\phiw{}\}\approx 2\qe{}$ do not vanish, so $\dot{\chi}_i\approx 0$ solve for the $u_i$ and terminate the algorithm with \emph{zero} dynamical d.o.f. In the nonlinear theory $\dot{\phi}_i\approx 0$ \emph{already} solve for the $u_i$, the $\chi_i$ are not induced, and a single \dof (two unconstrained Cauchy data) is \emph{activated}~\cite{Golovnev:2022rui}. Next, with $\lambda$ we additionally have $\phil{}\equiv\pil{}\approx 0$ and this time it is a trivial exercise to confirm that background linearisation \emph{changes nothing} (see~\cref{ConstraintAlgebra}). In this metaphor, we compare $\lambda$ in~\cref{MWE} with $\TLambda[1]{^\sigma_{\mu\nu}}$ in~\cref{DC}, the effect is to drag the nonlinear \dof down into the linear regime, and thereby solve a major problem in non-Riemannian gravity~\cite{Yo:2001sy,BeltranJimenez:2019hrm,Barker:2021oez,Kopczynski:1982,Cheng:1988zg,Dimakis:1989az,Chen:1998ad,Dimakis:1989ba,Blixt:2020ekl,Lemke:1990su,1999IJMPD...8..459Y,Hecht:1990wn}.

\begin{figure}[t]
  \center
  \includegraphics[width=\linewidth]{CombinedTheories.pdf}
	\caption{\label{ConstraintAlgebra} 
	Our solution to non-Riemannian strong coupling can easily be understood through the 1D minimal working example in~\cref{MWE}. Without our method, the $\chi_e$, $\chi_\omega$ constraints convert into a nonlinear torsion d.o.f. Accordingly, in actual nonlinear gravity~\cite{Yo:2001sy,Barker:2021oez}, constraints are converted into the $\B{_{\mu\nu}}$ field in~\cref{KalbRamondEffectiveTheory} (see~\cref{FullAnalysis}).
	}
\end{figure}

\section{Application to dark matter}


We can broadly distinguish two situations, according to whether or not the Proca field $\T[2]{_{\mu}}$ propagates at the relevant energy scales. The first case has been considered in~\cite{Ford:1989me,Golovnev:2008cf,Chiba:2008eh,Koivisto:2008xf,Kanno:2008gn,Emami:2016ldl,Armendariz-Picon:2004say,Koivisto:2007bp,Koivisto:2008xf,Tasinato:2014eka,DeFelice:2016yws,DeFelice:2016uil,BeltranJimenez:2016afo,deFelice:2017paw,Nakamura:2018oyy,Heisenberg:2020xak,Benisty:2021sul,deRham:2021efp,Hambye:2008bq,Hambye:2009fg,Arina:2009uq,Hisano:2010yh,Diaz-Cruz:2010czr,Lebedev:2011iq,Farzan:2012hh,Baek:2012se,Arcadi:2020jqf,Barman:2021qds,Bekenstein:1971hc,Heisenberg:2017hwb,Garcia-Saenz:2021uyv}. Whereas all these works used as a starting point an \emph{effective} Lagrangian with a structure similar to \cref{MasterLagrangian}, our theory \eqref{DC} can endow such models with a fundamental -- and gravitational -- origin.


Formula \eqref{Mass2} for the vector mass hints towards the second case, in which torsion does not propagate at the relevant energy scales. Namely, we already know that $\alp{5}<0$ and $\bet{2}>-1/3$. If $\bet{2}$ is not too close to the border value (\eg $\bet{2}\gtrsim0$) and $|\alp{5}|$ is at most of order one, we observe that $\mass[2]{}^2 \gtrsim \Planck^2$. Thus, one is tempted to conclude that parameter choices for which the Proca field does not propagate below the Planck scale are more generic. This is good news since then the energy scale at which scattering amplitudes involving the massive vector field violate perturbative unitarity also lies above the Planck scale. Consequently, introducing propagating torsion with the theory \eqref{DC} generically does not lower the cutoff scale, above which the theory ceases to be predictive. 

Integrating out the field $\T[2]{_{\mu}}$ at energies below $\mass[2]{}$, we get from \eq \eqref{MasterLagrangian} (with  $\T[3]{_{\mu}}\to 0$)
\begin{equation} \label{effectiveLagrangian}
	\mathcal{L}_{\text{T}} = - \frac{\Planck^2}{2}  \rR{} - \frac{1}{18 \mass[2]{}^2} \Si[2]{_{\mu}} \Si[2]{^{\mu}} + 	L_{\text{M}}\left(\rCon{}\right) \;.
\end{equation}
As a concrete example, we shall focus on fermionic fields in $L_{\text{M}}\left(\rCon{}\right)$ and first consider a generic $\Si[2]{_{\mu}} = -3 \sum_j \big(\zeta_V V^{(j)}_\mu + \zeta_A A^{(j)}_\mu\big)$, where $j$ sums species, $\zeta_V$ and $\zeta_A$ are real constants and $V_\mu = \bar{\Psi} \gamma_\mu \Psi$ and $A_\mu = \bar{\Psi} \gamma_5 \gamma_\mu \Psi$ represent the fermionic vector and axial currents, respectively. Following~\cite{Shaposhnikov:2020aen}, we shall take into account the field content of the Standard Model and add to it a singlet fermion $N$ of mass $M_N$. Such a scenario is strongly motivated since a small mixing with active neutrinos can generate neutrino masses via the seesaw mechanism (see~\cite{Boyarsky:2018tvu} for a review). For the present discussion, however, we shall not specify the precise nature of $N$. Independently of a possible connection to neutrinos, a singlet fermion is a dark matter candidate, but this option is only viable if a sufficient production of $N$ takes places in the early Universe. 

This can be achieved both by a propagating torsion field acting as mediator~\cite{Barman:2019mlj} and by the effective 4-fermion interaction arising from \eq \eqref{effectiveLagrangian} ~\cite{Shaposhnikov:2020aen}. We focus on the second case and consider the parameter choice $\zeta_V \gg 1$ and $\zeta_A \gg 1$, which can arise from a non-minimally coupled kinetic term of the fermion, $i/2 \bar{\Psi} (1+ 2 i \zeta_V + 2 i \zeta_A \gamma_5) \gamma^\mu \mathcal{D}_\mu \Psi + \text{h.c.}$ (see~\cite{Shaposhnikov:2020frq,Karananas:2021zkl,Rigouzzo:2023sbb}).  Then the produced abundance of fermions is~\cite{Shaposhnikov:2020aen}
\begin{equation}
	\frac{\Omega_N}{\Omega_{DM}}\simeq 1.7 \, C \, \frac{M_P}{\mass[2]{}} \left( \frac{M_N}{10~\text{keV}}\right) \, \left( \frac{T_\mathrm{prod}}{\mass[2]{}} \right)^3 \;.
	\label{darkMatterAbundance}
\end{equation}
Here $T_\mathrm{prod}$ is the highest temperature at which fermions are generated, \ie generically the temperature of the hot Big Bang, and $C = \left(15(\zeta_V^2 - \zeta_A^2)^2 + 7(\zeta_V + \zeta_A)^4 + 8 (\zeta_V - \zeta_A)^4\right)$ for a Dirac fermion $N$. Choosing $C$ sufficiently large, producing all of dark matter, $\Omega_N=\Omega_{DM}$, can be achieved for a wide range of masses $M_N$ down to $\text{few}~\text{keV}$, where even lighter fermions are excluded due to well-known observational bounds on warm dark matter~\cite{Garzilli:2018jqh, Garzilli:2019qki}. The resulting mass hierarchy typically is $\mass[2]{}\gg M_P \gg T_\mathrm{prod}$. Also $\mass[2]{}< M_P$ is possible, as long as $T_\mathrm{prod}<\mass[2]{}$ so that $\T[2]{_{\mu}}$ does not propagate.

\section{Conclusions} 
Several long-standing problems related to gravity have remained unsolved. Since many of the proposed solutions call for the introduction of new particles, one naturally wonders if GR can provide them and indeed it is easy to generate additional propagating degrees of freedom, \eg from higher powers of curvature terms. However, such models are generically plagued by problems due to ghost instabilities~\cite{Stelle:1977ry,Neville:1978bk,Neville:1979rb,Sezgin:1979zf,Hayashi:1979wj, Hayashi:1980qp} and strong coupling~\cite{Hecht:1996np,Yo:1999ex,Yo:2001sy}. In this paper we addressed the notorious problem of \emph{linear constraints converting into nonlinear ghosts} in the torsion sector~\cite{Yo:2001sy,Hayashi:1980qp,Blagojevic:1983zz,Barker:2021oez,Blagojevic:2013dea,Blagojevic:2013taa,Aoki:2020rae,BeltranJimenez:2019hrm,Hecht:1996np,Yo:1999ex,Yo:2001sy,Hecht:1991jh,Magueijo:2008sx,Afshordi:2006ad,Dimakis:1989az,Dimakis:1989ba,Lemke:1990su,Hecht:1990wn,Delhom:2022vae,Velo:1969txo,Aragone:1971kh,Krasnov:2021zen,Chen:1998ad,Cheng:1988zg,Chen:1998ad,Blixt:2020ekl}. Direct canonical methods in~\cref{DC,MWE} suggest there is in fact a loophole in \emph{vector torsion}. This is a significant new example of propagating vectors arising from a curvature-squared invariant, for which consistency can be proven at the full nonlinear level.

Our theory \eqref{DC} endows effective Einstein-Proca models with a fundamental gravitational origin. This leads to significant constraints on the parameter space even in situations in which torsion does not propagate, and one can say that the remaining freedom is `just right'. On the one hand, it is possible to address some of the unresolved puzzles in cosmology with approaches that do not exist in the metric formulation of GR. As an example, we discussed a novel mechanism for producing fermionic dark matter~\cite{Shaposhnikov:2020aen}. On the other hand, theories such as the one that we have constructed generically feature significantly fewer free parameters as compared to effective Einstein-Proca models~\cite{Ford:1989me,Golovnev:2008cf,Chiba:2008eh,Koivisto:2008xf,Kanno:2008gn,Emami:2016ldl,Armendariz-Picon:2004say,Koivisto:2007bp,Koivisto:2008xf,Tasinato:2014eka,DeFelice:2016yws,DeFelice:2016uil,BeltranJimenez:2016afo,deFelice:2017paw,Nakamura:2018oyy,Heisenberg:2020xak,Benisty:2021sul,deRham:2021efp,Hambye:2008bq,Hambye:2009fg,Arina:2009uq,Hisano:2010yh,Diaz-Cruz:2010czr,Lebedev:2011iq,Farzan:2012hh,Baek:2012se,Arcadi:2020jqf,Barman:2021qds,Bekenstein:1971hc,Heisenberg:2017hwb,Garcia-Saenz:2021uyv} or proposals~\cite{Freidel:2005sn,Bauer:2008zj,Poplawski:2011xf,Diakonov:2011fs,Khriplovich:2012xg,Magueijo:2012ug,Khriplovich:2013tqa,Markkanen:2017tun,Carrilho:2018ffi,Enckell:2018hmo,Rasanen:2018fom,Rubio:2019ypq,Shaposhnikov:2020geh,Karananas:2020qkp,Langvik:2020nrs,Shaposhnikov:2020gts,Mikura:2020qhc,Shaposhnikov:2020aen,Kubota:2020ehu,Enckell:2020lvn,Iosifidis:2021iuw,Racioppi:2021ynx,Cheong:2021kyc,Dioguardi:2021fmr,Piani:2022gon,Dux:2022kuk,Rasanen:2022ijc,Gialamas:2023emn,Gialamas:2023flv,Piani:2023aof,Poisson:2023tja,Rigouzzo:2023sbb} with non-propagating torsion, and hence are more predictive.



Several directions for future research emerge from our findings. First, it remains to be determined if our theory is special or if classes of models with analogous properties exist in Einstein-Cartan gravity and other formulations of GR. Second, our model can constrain the parameter space in inflationary models that suffer from a loss of predictivity due to numerous unknown coupling constants (see \eg~\cite{Bauer:2008zj,Rasanen:2018ihz, Raatikainen:2019qey, Langvik:2020nrs, Shaposhnikov:2020gts,Rigouzzo:2022yan}). Third, it would be interesting to explore further observables consequences such as in gravitational waves~\cite{Battista:2021rlh}. Finally, consistent gravitational theories with curvature-squared terms may provide new approaches for the ultimate challenge of UV-completing GR~\cite{Lin:2019ugq}.

\begin{acknowledgments}
This work was performed using resources provided by the Cambridge Service for Data Driven Discovery (CSD3) operated by the University of Cambridge Research Computing Service (\href{www.csd3.cam.ac.uk}{www.csd3.cam.ac.uk}), provided by Dell EMC and Intel using Tier-2 funding from the Engineering and Physical Sciences Research Council (capital grant EP/T022159/1), and DiRAC funding from the Science and Technology Facilities Council (\href{www.dirac.ac.uk}{www.dirac.ac.uk}).

We thank Jaakko Annala, Daniel Blixt, Francisco José Maldonado Torralba, Carlo Marzo, Roberto Percacci, Syksy Räsänen, Claire Rigouzzo, Inar Timiryasov and Tom Złośnik for useful comments. W.B. is grateful for the kind hospitality of Leiden University and the Lorentz Institute, and the support of Girton College, Cambridge. S.Z.~acknowledges support of the Fonds de la Recherche Scientifique -- FNRS.
\end{acknowledgments}



\appendix

\begin{figure}[t]
  \center
  \includegraphics[width=\linewidth]{PhaseSpace.pdf}
	\caption{\label{PhaseSpace} 
	Strong coupling splits the perturbative vacuum. Left: the linear approximation to~\cref{ToyModel} erroneously anticipates the stable vacuum. Right: in the bulk of the true nonlinear phase space this vacuum is torn apart by a high-energy separatrix.
	}
\end{figure}

\begin{figure*}[h]
  \center
	\vspace{20pt}
  \includegraphics[width=\linewidth]{SupplementalCombinedTheories.pdf}
	\caption{\label{FullAnalysisFigure} 
	On-shell constraint algebra of $\LPV{(2)\rightleftarrows(3)}$ in~\eqref{DC}, with and without our multipliers, with the logical flow of the \emph{linear} Dirac algorithm superimposed. This figure supplements~\cref{ConstraintAlgebra} for completeness. Without our method, the induced secondary constraints $\tensor*{\chi}{_e^{1^-}}$, $\tensor*{\chi}{_\omega^{0^-}}$, $\tensor*{\chi}{_\omega^{2^-}}$ of the theory~\cref{DC} would collectively lose six (canonical) \dof in nonlinear gravity~\cite{Yo:2001sy,Barker:2021oez}, activating half of the $\B{_{\mu\nu}}$ field in~\cref{KalbRamond}. If $\bet{2}\to 0$ had additionally been enforced, the whole of $\B{_{\mu\nu}}$ would have become activated as per the analysis in~\eqref{KalbRamondEffectiveTheory}. The reason we do not show the nonlinear algebra in this case is because the loss of secondaries occurs \emph{among} specific components of given spin sectors, so the convenient spin-parity decomposition of the algebra falls apart. The full nonlinear constraint chains are listed in~\cite{HamiltonianAnalysis}.
	}
\end{figure*}

\section{Conventions}\label{Conventions}

In this appendix we quote our conventions. The metric $g_{\mu\nu}$ has signature $\left(+,-,-,-\right)$ and the covariant derivative of a covector $A_\nu$ is $\rcD{_\mu} A_\nu = \partial_\mu A_\nu - \Gamma^\alpha_{~~ \mu\nu} A_\alpha$ with an a priori independent affine connection $\Gamma^\alpha_{~~\beta \gamma}$ and torsion $\tensor{T}{^\mu_{\nu\sigma}}\equiv -2\rcCon{^\mu_{[\nu\sigma]}}\equiv\frac{4}{3}\T[1]{^\mu_{[\nu\sigma]}}+\frac{2}{3}\tensor*{\delta}{^\mu_{[\nu}}\T[2]{_{\sigma]}}+\tensor{\epsilon}{^\mu_{\nu\sigma\lambda}}\T[3]{^\lambda}$ decomposed into tensor $\T[1]{^\lambda_{\mu\nu}}\equiv\frac{1}{2}\tensor{T}{^\lambda_{\mu\nu}}+\tfrac{1}{2}\tensor{T}{_\mu^\lambda_\nu}+\tfrac{1}{6}\tensor{g}{_{\mu\nu}}\T[2]{^\lambda}-\tfrac{1}{6}\tensor*{\delta}{^\lambda_\nu}\T[2]{_\mu}+\tfrac{1}{3}\tensor*{\delta}{^\lambda_\mu}\T[2]{_\nu}$, vector $\T[2]{_\mu}\equiv\tensor{T}{^\lambda_{\lambda\mu}}$ and pseudovector $\T[3]{_\mu}\equiv\tfrac{1}{6}\tensor{\epsilon}{_{\mu\nu\sigma\lambda}}\tensor{T}{^{\nu\sigma\lambda}}$ parts, with nonmetricity $Q_{\alpha\mu\nu}\equiv\nabla_\alpha g_{\mu\nu}$, where anti-symmetrization on some tensor $T_{[\mu\nu]}$ is defined as $T_{[\mu\nu]} \equiv \frac{1}{2}(T_{\mu\nu}-T_{\nu\mu})$. The curvature has contractions $\rcR{}\equiv\rcR{^\lambda_\lambda}\equiv\rcR{^{\mu\nu}_{\mu\nu}}$, with $\rcR{_\lambda^\sigma_{\mu\nu}}\equiv 2\PD{_{[\nu}}\rcCon{^\sigma_{\mu]\lambda}}+2\rcCon{^\sigma_{[\nu|\kappa}}\rcCon{^\kappa_{|\mu]\lambda}}$, and the Riemann curvature $\rR{_\lambda^\sigma_{\mu\nu}}$ is analogously defined in terms of the metric-compatible and torsion-free Levi-Civita connection $\rCon{^\sigma_{\mu\nu}}$, with covariant derivative $\rD{_\mu}$. We also use the tetrad $\tensor{e}{_i^\mu}\tensor{e}{_j^\nu}\tensor{\eta}{^{ij}}\equiv\tensor{g}{^{\mu\nu}}$, with $e\equiv\det\tensor{e}{^i_\mu}\equiv\sqrt{-g}\equiv\sqrt{-\det\tensor{g}{_{\mu\nu}}}$, and independent spin connection $\tensor{\omega}{^{ij}_\mu}$. For our gamma matrices, we use the convention $\left\{\gamma_i, \gamma_j \right\} \equiv - 2 \eta_{ij}$, $\gamma_5 \equiv i \gamma_0 \gamma_1 \gamma_2 \gamma_3$, with Greek indices on gamma matrices implying contraction with an (inverse) tetrad field, and in flat space with Cartesian coordinates the Levi-Civita tensor has the same components as the symbol, $\epsilon_{0123}\equiv 1$.

\section{Toy-model of strong coupling}\label{SCIntro}


\vspace{10pt}

In this appendix we clarify what we mean by \emph{strong coupling}, and why it is commonly diagnosed in the Hamiltonian framework. Consider the double-oscillator toy model (TM) of \dof $q_1$ and $q_2$ in one dimension
\begin{equation}\label{ToyModel}
	L_{\text{TM}}=\frac{1}{2}\dot{q}_1^2-\frac{1}{2}q_1^2+\frac{1}{2}q_1^2\dot{q}_2^2-\frac{1}{2}q_2^2 \;.
\end{equation}
The perturbative `approximation' to~\cref{ToyModel} is $L_{\text{TM}}=\dot{q}_1^2/2-q_1^2/2-q_2^2/2+\mathcal{O}\left(q^4\right)$, for which the second oscillator drops out algebraically as $q_2\approx 0$, leaving the Hamiltonian phase subspace $\left(q_1,p_1\equiv\dot{q}_1\right)$ in which the first oscillator orbits the perfectly healthy vacuum at $q_1\approx p_1\approx 0$ with Hamilton's equations $\left(\dot{q}_1\approx p_1,\dot{p}_1\approx-q_1\right)$. However, the true physics always lives in the space $\left(q_1,p_1,q_1,p_2\right)$, in which the Hamilton equation $\dot{p}_1\approx -q_1+p_2^2/q_1^2$ is \emph{inherently non-perturbative}. As illustrated in~\cref{PhaseSpace}, the effect is absolutely fatal to the perturbative vacuum. Generalising to field theory, strongly-coupled \dof, analogous to $q_2$, reveal themselves via structural changes to the Hamiltonian constraint structure when passing from linearised to nonlinear gravity theories. In such cases it is not clear how the background of the linearisation can then be a viable spacetime, even though it may perfectly well be an exact solution to the nonlinear field equations. The sick kinetic structure of~\eqref{ToyModel} may be compared to that in~\eqref{KalbRamondEffectiveTheory}.

It is important to point out that the splitting of the vacuum in~\cref{PhaseSpace} is an effect which \emph{need not} be associated with ghost modes. Indeed, in~\cref{ToyModel} we are careful to preserve the sign of the kinetic term. In our practical, field theoretic example in~\cref{KalbRamondEffectiveTheory} however, it may turn out to be harder to bound the Hamiltonian due to the multiple components of $\B{_{\mu\nu}}$. 

\section{Full nonlinear Hamiltonian analysis}\label{FullAnalysis}

In this appendix, we quote the full, nonlinear Hamiltonian analysis of $\LPV{(2)\rightleftarrows(3)}$ corresponding to $\alp{5}>0$ w.l.o.g. The full Dirac algorithm is shown schematically in~\cref{FullAnalysisFigure}, using precisely the same notation as in~\cref{ConstraintAlgebra}. Note the appearance of tertiary $\zeta$ and quaternary $\theta$ constraints. As before, the full nonlinear constraint chains are found in~\cite{HamiltonianAnalysis}. The key observation is that the size of the algebra (i.e. the number of induced non-primary constraints) does not change in passing from the linear to the nonlinear theories (though we do not exclude this effect at other points which punctuate the bulk of the phase space). This is consistent with the \emph{manifestly} healthy and dynamically equivalent Einstein-Proca model.

\bibliography{Manuscript}

\end{document}
